\RequirePackage{fix-cm}
\newif\ifreviewfallback
\IfFileExists{elsarticle.cls}{
  \reviewfallbackfalse
  \documentclass[5p,times,authoryear]{elsarticle}
}{
  \reviewfallbacktrue
  \documentclass[10pt,a4paper,twocolumn]{article}
}
\ifreviewfallback\usepackage{review-compat}\fi
\input{corpus_counts.tex}


\usepackage{graphicx}
\usepackage{amsmath}            % Equation environments and cross-references
\usepackage{epstopdf}           % EPS support when compiling with pdfLaTeX
% Avoid forced two-column balancing during peer-review drafts.
\usepackage{amssymb}
\usepackage{xurl}
\usepackage{adjustbox}
\usepackage{tikz}
\usetikzlibrary{positioning,arrows.meta}
\usepackage{booktabs} % For \toprule, \midrule, \bottomrule in tables
\usepackage{array}
\usepackage{tabularx}
\usepackage{longtable}
\usepackage{placeins}
\usepackage[T1]{fontenc}
\usepackage[utf8]{inputenc}
\usepackage{microtype} % Improved kerning and line breaking
\usepackage{hyperref}
\hypersetup{hidelinks,pdftitle={Architectural Transitions in Medical Image Segmentation},pdfauthor={Nand Kumar Yadav and colleagues}}
\usepackage{multirow} % For \multirow in tables
\usepackage{xcolor} % Color support for diagrams and tables

\usepackage[main=english]{babel} % English prose; bibliography strings are set by the bundled .bst
\AtBeginDocument{\selectlanguage{english}}
% Improve double-column figure placement
\setcounter{dbltopnumber}{2}
\renewcommand{\dbltopfraction}{0.95}
\renewcommand{\textfraction}{0.05}
\renewcommand{\dblfloatpagefraction}{0.80}
% Landscape tables
\usepackage[figuresright]{rotating}

% Multi-panel figures are constructed from minipages (no deprecated subfigure package).
\setlength{\emergencystretch}{1.5em} % reduces overfull lines in two-column text

% declarations for front matter

\usepackage{helvet}
\usepackage{courier}

\begin{document}

\begin{frontmatter}


\title{A Taxonomy of Architectural Transitions in Medical Image Segmentation: A Framework for Bottleneck Redistribution}


\author[Second]{Nand Kumar Yadav\corref{cor1}}
\ead{nand.yadav@usd.edu}


\author[First]{Smriti Regmi}
\ead{smriti.regmi@coyotes.usd.edu}

\author[First]{Moinak Bose}
\ead{moinak.bose12@gmail.com}

\author[Second]{N. Reed Alexander}
\ead{reed.alexander@usd.edu}


\author[First]{Debesh Jha}
\ead{debesh.jha@usd.edu}
\author[Second]{William CW Chen}
\ead{william.chen@usd.edu}


\cortext[cor1]{Corresponding author.}


\address[First]{Department of Computer Science, University of South Dakota, USA}
\address[Second]{Department of Biomedical and Translational Sciences, Sanford School of Medicine, University of South Dakota, Vermillion, SD, USA}


\begin{abstract}
Medical image segmentation has advanced through convolutional encoder--decoder networks, transformers, state-space models, and promptable foundation models. Reported gains, however, often combine architectural changes with differences in initialization, input geometry, training, and evaluation. We present a structured narrative synthesis of \CorpusTotal{} named medical segmentation method families, documented in an accompanying method inventory. Four analytical objectives organize the review: local spatial precision, global contextual reasoning, computational efficiency, and cross-domain generalization. We examine six classes of architectural intervention through a bounded redistribution hypothesis: under a specified protocol, which limitation is relieved and which prespecified constraint, if any, becomes limiting? The framework allows simultaneous improvements and requires a predicted constraint shift to be measured rather than inferred from unmatched benchmarks. Source-specific evidence records distinguish original evaluations, independent reimplementations, secondary reports, and unresolved values. Across three worked contrasts, the intended benefit is better documented than the predicted constraint shift, for which matched spatial, computational, or generalization tests remain incomplete. Mechanism-linked diagnostic tests connect these evidence gaps to five research directions: volumetric state propagation, medical pretraining, uncertainty-conditioned fusion, frequency-aware boundary learning, and multi-objective evaluation. The resulting synthesis provides an auditable basis for assessing architectural claims without implying a universal ranking or an unavoidable trade-off.
\end{abstract}

\begin{keyword}
Biomedical Image Segmentation; U-Net; Transformer; Swin Transformer; Mamba; Segment Anything Model (SAM); Deep Learning; Medical Imaging.
\end{keyword}

\end{frontmatter}


\section{Introduction}

While many architectures for medical image segmentation have been proposed and reported to achieve excellent results, explaining these results in terms of information-processing mechanisms remains challenging. A higher Dice score may reflect a different cohort split, input resolution, training strategy, or post-processing rather than architecture alone. We therefore ask three questions that cannot be resolved by a leaderboard: (i) What capability does an intervention add? (ii) Under which conditions is it useful? (iii) Which additional constraint, if any, becomes limiting?

Four design goals are frequently mentioned in the literature for medical image segmentation: local spatial precision, global contextual reasoning, computational efficiency, and cross-domain generalization. Each of these goals places a distinct demand on a segmentation system. For example, convolutional operators provide strong local spatial precision and translation bias, which are beneficial for preserving boundaries, small structures, and low-contrast regions~\citep{ronneberger2015unet,oktay2018attentionunet,isensee2021nnu}. Self-attention allows for direct interaction between distant regions, but the cost of self-attention increases rapidly with the number of tokens~\citep{dosovitskiy2021vit,chen2021transunet,hatamizadeh2022unetr}. Selective state-space models provide another form of long-range propagation, but vision models must first decide how to order a two- or three-dimensional anatomy for sequential processing~\citep{gu2023mamba}. Promptable models such as SAM provide large-scale pretrained representations and can be flexibly used with user guidance, but their performance on medical images still depends on domain adaptation, the quality of the prompts, and consistency across slices or volumes~\citep{kirillov2023sam,ma2024medsam}.

Each architectural family addresses particular information-processing demands and can therefore coexist with others within a single model. Our analysis uses a \emph{framework for bottleneck redistribution} to examine these combinations. For each architectural intervention, we identify its intended benefit and specify the conditions under which a predicted constraint shift could be tested. Each hypothesis is bounded by a task, an evaluation protocol, a target benefit, and prespecified constraints. An improvement in one metric does not establish that a cost has shifted elsewhere; when multiple objectives improve, the predicted constraint must still be evaluated with an appropriate stress test.

Reviews have already extensively discussed the sensitivity of medical image segmentation systems to the chosen architecture, data quality and training conditions (e.g. \citep{azad2024unet_verified,tajbakhsh2020embracing}). The current synthesis focuses on a more specific question, i.e. whether the proposed intervention achieves its intended benefit and whether it results in the expected shift in a declared constraint. For each of the six transitions described in Sections~\ref{subsec:failure_modes}--\ref{subsubsec:redistribution} and summarized in Table~\ref{tab:failure_matrix}, we identify the targeted limitation, the proposed mechanism and the candidate constraint to be tested.

Bottleneck redistribution should not be interpreted as an unavoidable deterioration in a headline metric. An intervention may improve several outcomes simultaneously, while a residual limitation becomes apparent only under a targeted stress test. Windowed attention illustrates the distinction: restricting direct token interactions reduces the computational burden of full attention, but can make long-range communication depend on shifted windows, depth, or additional connections. Whether this restriction produces task-relevant errors is an empirical question. A controlled comparison that varies cross-window connectivity while holding the remaining conditions fixed would support the explanation if cases requiring distant information respond selectively; absence of that pattern would weaken it.

The six transitions that are analyzed in this review-level synthesis are direct-to-adaptive skip transfer, CNN--transformer hybridization, hierarchical and efficient attention, SSM/Mamba state propagation, medical adaptation of promptable foundation models, and conditional routing. Results are presented separately from the review-level synthesis, and numerical values are reported in the protocol in which they were obtained, i.e. they are not treated as interchangeable.

An inventory of the \CorpusTotal{} medical segmentation method families counted in the retrospective synthesis (including CNN/U-Net, transformer, SSM/Mamba, SAM-based, and adaptive-routing architectures) is presented in the accompanying inventory. This inventory aims to cover the mechanisms used for representation building, information exchange across scales, spatial detail recovery, computational resources, and adaptation to changes in task, modality, or domain, rather than at an exhaustive list of all architectures.
The review makes four main contributions:
\begin{itemize}
    \item an information-flow taxonomy that places CNN, transformer, SSM/Mamba, promptable foundation-model, and adaptive-routing methods within a common four-objective design space;
    \item a failure analysis that links recurring weaknesses to plausible architectural causes while separating reported evidence from review-level interpretation;
    \item a structured comparison of six architectural transitions, supported by source-specific evidence records that state what was verified and what remains unresolved;
    \item a research roadmap that turns unresolved limitations into testable directions for three-dimensional state-space modeling, medical-native pretraining, adaptive feature fusion, frequency-aware boundary reconstruction, and multi-objective evaluation.
\end{itemize}

Figure~\ref{fig:taxonomy} gives the mechanism taxonomy and Figure~\ref{fig:redistribution_map} visualizes the six bottleneck transitions. The five roadmap hypotheses are developed cumulatively from these architecture-specific bottlenecks, the mechanism-linked failure matrix, and the protocol-audited evidence analysis. They are therefore downstream hypotheses of the taxonomy rather than a detached future-work list.

Table~\ref{tab:survey_comparison} places this review alongside representative surveys of deep learning, CNN--transformer segmentation, SSM/Mamba models, and biomedical segmentation foundation models. The comparison is intended to clarify scope rather than rank review quality. Several earlier surveys provide deeper coverage of individual model families. The present review instead brings these families together within one framework and examines how their strengths and limitations relate across a common set of design objectives.


\begin{table*}[t]
\centering
\caption{Positioning against representative reviews available within the coverage window. The comparison identifies emphasis, not a ranking of review quality. The present contribution is the connection between a bounded mechanism hypothesis, a source-specific evidence record, and a discriminating test.}
\label{tab:survey_comparison}
\small\renewcommand{\arraystretch}{1.15}\setlength{\tabcolsep}{5pt}
\begin{tabularx}{\textwidth}{p{3.25cm}p{4.6cm}X}
\toprule
\textbf{Review} & \textbf{Main emphasis} & \textbf{Relationship to the present synthesis} \\
\midrule
\citep{litjens2017survey} & Broad deep learning methods in medical image analysis. & Establishes the wider methodological setting and clinical tasks. \\
\citep{shamshad2023transformers} & Transformer methods in medical imaging. & Provides deeper attention-family coverage; our analysis connects attention restrictions to mechanisms in other families. \\
\citep{azad2024unet_verified} & U-Net variants and their segmentation evaluations. & Provides detailed encoder--decoder analysis and medical reimplementations; we use those evaluations with their source protocols. \\
\citep{shu2026efficient_verified} & Efficient segmentation, including Mamba, foundation-model adaptation, and lightweight architectures. & Establishes substantial contemporary cross-family overlap; breadth and efficiency alone do not distinguish this review. \\
\citep{moglia2025generalist_verified} & Generalist segmentation models and their relationship to specialist systems. & Provides a focused treatment of foundation/generalist models; our framework also follows reconstruction, serialization, and routing constraints. \\
This review & Architectural interventions, evidence provenance, and diagnostic tests across families. & States a bounded hypothesis for each transition, distinguishes incompatible measurements, and links remaining questions to controlled contrasts. \\
\bottomrule
\end{tabularx}
\end{table*}


The scope table should therefore be read as a map of emphasis, not as a quality ranking. The present review contributes a shared LP--GC--CE--CG vocabulary, the redistribution framework as a bounded diagnostic approach, a mechanism-linked failure matrix (Table~\ref{tab:failure_matrix}), and a research roadmap (Table~\ref{tab:research_roadmap}).
\begin{figure*}[t]
\centering
\begingroup
\definecolor{ReviewCNN}{HTML}{0646B5}%
\definecolor{ReviewCNNFill}{HTML}{E8F0FF}%
\definecolor{ReviewAttention}{HTML}{6B55C9}%
\definecolor{ReviewAttentionFill}{HTML}{E5DEFD}%
\definecolor{ReviewSSM}{HTML}{4A1FA0}%
\definecolor{ReviewSSMFill}{HTML}{EEE3FF}%
\definecolor{ReviewSAM}{HTML}{1F5FBF}%
\definecolor{ReviewSAMFill}{HTML}{EEF5FF}%
\definecolor{ReviewFusion}{HTML}{2E7D32}%
\definecolor{ReviewFusionFill}{HTML}{D9F0D5}%
\definecolor{ReviewLine}{HTML}{555555}%

\begin{tikzpicture}[
    every node/.style={
        font=\sffamily\fontsize{9}{10.2}\selectfont,
        text=black
    },
    arr/.style={
        -{Latex[length=1.5mm]},
        line width=0.6pt,
        draw=ReviewLine
    },
    family/.style={
        rounded corners=2pt,
        line width=0.6pt,
        text width=37mm,
        minimum height=17mm,
        align=center,
        inner sep=1mm
    },
    route/.style={
        draw=ReviewFusion,
        fill=ReviewFusionFill,
        rounded corners=2pt,
        line width=0.6pt,
        text width=166mm,
        minimum height=6.5mm,
        align=center,
        inner sep=1mm
    }
]

\node[family,draw=ReviewCNN,fill=ReviewCNNFill]
    (cnn) at (0,0)
    {\textbf{Convolutional}\\
     \textbf{encoder--decoder}\\
     Local features and skips\\
     Spatial reconstruction};

\node[family,draw=ReviewAttention,fill=ReviewAttentionFill]
    (attn) at (43mm,0)
    {\textbf{Transformer and}\\
     \textbf{attention hybrids}\\
     Content-dependent\\
     token interactions};

\node[family,draw=ReviewSSM,fill=ReviewSSMFill]
    (ssm) at (86mm,0)
    {\textbf{Selective}\\
     \textbf{state-space model}\\
     Compressed recurrent state\\
     Directional scans};

\node[family,draw=ReviewSAM,fill=ReviewSAMFill]
    (sam) at (129mm,0)
    {\textbf{Promptable}\\
     \textbf{foundation model}\\
     Pretraining and prompts\\
     Medical / 3D adaptation};

\node[route] (routing) at (64.5mm,-17mm)
    {\textbf{Cross-cutting routing and fusion:}
     conditional selection of features, scales, experts or computation};

\foreach \n in {cnn,attn,ssm,sam}{
    \draw[arr] (\n.south) -- (\n.south |- routing.north);
}

\end{tikzpicture}
\endgroup
\caption{Information-flow taxonomy. Families can coexist; routing and fusion
operate across them. The layout implies no chronology or ranking. Each method
is counted under one primary family while retaining its hybrid description. Drawing code and layout were prepared with assistance from ChatGPT (OpenAI; historical model version not recorded).}
\label{fig:taxonomy}
\end{figure*}

\begin{figure*}[t]
\centering
\begingroup
\definecolor{ReviewCNN}{HTML}{0646B5}
\definecolor{ReviewCNNFill}{HTML}{E8F0FF}
\definecolor{ReviewAttention}{HTML}{6B55C9}
\definecolor{ReviewAttentionFill}{HTML}{E5DEFD}
\definecolor{ReviewSSM}{HTML}{4A1FA0}
\definecolor{ReviewSSMFill}{HTML}{EEE3FF}
\definecolor{ReviewSAM}{HTML}{1F5FBF}
\definecolor{ReviewSAMFill}{HTML}{EEF5FF}
\definecolor{ReviewFusion}{HTML}{2E7D32}
\definecolor{ReviewFusionFill}{HTML}{D9F0D5}
\definecolor{ReviewConcern}{HTML}{C00000}
\definecolor{ReviewConcernFill}{HTML}{FFF0F0}
\definecolor{ReviewLine}{HTML}{555555}
\begin{tikzpicture}[
    every node/.style={font=\sffamily\fontsize{9}{10.8}\selectfont,text=black},
    arr/.style={-{Latex[length=2mm]},line width=0.7pt,draw=ReviewLine},
    card/.style={rounded corners=3pt,line width=0.7pt,text width=76mm,
        minimum height=22mm,align=flush left,inner sep=2mm}]
\node[card,draw=ReviewCNN,fill=ReviewCNNFill] at (44mm,0)
    {\textbf{1. Direct skips $\rightarrow$ selective fusion}\\Target: retain useful fine-scale features.\\Constraint to test: alignment, noise transfer, routing cost.};
\node[card,draw=ReviewAttention,fill=ReviewAttentionFill] at (130mm,0)
    {\textbf{2. Local encoding $\rightarrow$ explicit attention}\\Target: connect distant anatomical evidence.\\Constraint to test: token cost, resolution, initialization.};
\node[card,draw=ReviewAttention,fill=ReviewAttentionFill] at (44mm,-28mm)
    {\textbf{3. Full $\rightarrow$ restricted attention}\\Target: make high-resolution interaction tractable.\\Constraint to test: communication across windows or sparse edges.};
\node[card,draw=ReviewSSM,fill=ReviewSSMFill] at (130mm,-28mm)
    {\textbf{4. Attention $\rightarrow$ state propagation}\\Target: avoid a dense token-interaction matrix.\\Constraint to test: ordering, state compression, spatial coherence.};
\node[card,draw=ReviewSAM,fill=ReviewSAMFill] at (44mm,-56mm)
    {\textbf{5. Generic $\rightarrow$ medical foundation model}\\Target: improve task and domain alignment.\\Constraint to test: prompting, adaptation cost, retained transfer.};
\node[card,draw=ReviewFusion,fill=ReviewFusionFill] at (130mm,-56mm)
    {\textbf{6. Fixed $\rightarrow$ conditional allocation}\\Target: allocate capacity to the current input.\\Constraint to test: calibration, utilization, execution overhead.};
\end{tikzpicture}
\endgroup
\caption{Six candidate architectural transitions. Each card separates the intended benefit from constraints that must be specified and measured in a controlled test. A listed constraint is a hypothesis, not an inevitable cost. Joint improvement can contradict a particular redistribution hypothesis within its declared scope. Drawing code and layout were prepared with assistance from ChatGPT (OpenAI; historical model version not recorded).}
\label{fig:redistribution_map}
\end{figure*}

\section{Review Design and Analytical Framework}
\label{sec:method}

Medical image segmentation is more than mapping an image to a pixel-wise label map. A segmentation model must preserve fine spatial detail, capture long-range anatomical relationships, operate within available computational resources, and remain reliable across scanners, institutions, patient populations, and imaging protocols. Improving one objective may change the constraints governing another, although simultaneous gains are possible. Architectural choices therefore need to be assessed within specified data, compute, and deployment conditions.

For a representative two-dimensional, $K$-class setting, let $X \in \mathbb{R}^{H \times W \times C}$ denote an input image and $Y \in \{1,\ldots,K\}^{H \times W}$ its corresponding label map. A model with parameters $\theta$ defines

\[
f_{\theta}: X \rightarrow Y.
\]

Architectural decisions about the spatial relations a model can represent, the information retained as resolution changes, and the required computation are grouped into three functions:

\begin{itemize}
    \item \textbf{Local feature extraction} – capturing local features such as edges, textures, and small structures.
    \item \textbf{Global context modeling} – capturing long-range anatomical and semantic relationships.
    \item \textbf{Feature fusion} – fusing information from different scales and representations during reconstruction.
\end{itemize}

There are different architectural families for these functions. CNNs use spatially local filters and hierarchical feature extraction, giving them strong local inductive bias. Transformers use self-attention to model long-range interactions between distant regions. State-space models propagate information through latent states and provide an alternative form of long-range modeling with lower asymptotic complexity than full attention. Foundation models add large-scale pretraining and prompt-based adaptation, shifting part of the problem from architecture design to representation transfer and task adaptation.

This also means that different architectural choices lead to different constraints. For example, more contextual interaction can lead to higher memory and computation requirements, as well as the need for more training data. In order to recover fine boundaries, one typically needs to keep high-resolution features throughout the network, and the decoder must be sufficiently large. Finally, a network that has been strongly adapted to a certain dataset can have very high accuracy on that dataset, but its performance can be very unstable when applied to other domains.

These questions will be used as a basis for the comparison of the different architectural families in the following sections instead of solely reporting the accuracy on the different benchmarks. For each family, we discuss how the representations are built, how anatomical context is handled, how information is propagated between different scales, how the computation is used and how the models behave on different tasks and in different domains.
\subsection{Analytical Constructs and Operational Evidence}
\label{subsec:constructs}

The four objectives defined in the previous section are used in the analysis of architectural behavior in Section 4. Note that each of the four objectives is not measured by a single metric. For example, local spatial precision is not measured by DSC alone and global contextual reasoning is not measured by HD95 alone. Rather, each metric captures some aspect of the objective in question.

Evidence for claims regarding the objectives is most convincing when it stresses the capability of interest. Thus, boundary and surface metrics are most relevant for local spatial accuracy, long-range interaction tests, inter-slice consistency, and context-specific ablations are most relevant for global reasoning, and computational efficiency is most relevantly assessed using matched implementation and input settings. Cross-domain generalization is most relevantly assessed on held-out domains that are defined using the same task and label definition.

Table~\ref{tab:construct_definitions} lists each objective with its corresponding interpretation and evidence that supports it. No single metric is used to represent an objective. Rather, the conclusions are based on measurements and stress tests that are most relevant to the claim being made.

\begin{table*}[t]
\centering
\caption{Operational interpretation of the four analytical objectives. The listed measurements are admissible evidence or stress tests, not interchangeable definitions of the construct. LP and GC require mechanism-specific diagnostics for strong claims; CE and CG also require matched accounting and task definitions, respectively.}
\label{tab:construct_definitions}
\footnotesize
\renewcommand{\arraystretch}{1.05}
\setlength{\tabcolsep}{3.0pt}

\begin{tabularx}{\textwidth}{>{\raggedright\arraybackslash}p{2.0cm}>{\raggedright\arraybackslash}p{4.1cm}>{\raggedright\arraybackslash}p{5.8cm}>{\raggedright\arraybackslash}X}
\toprule
\textbf{Construct} & \textbf{Analytical meaning} & \textbf{Admissible operational evidence} & \textbf{Not equivalent to} \\
\midrule
Local spatial precision (LP) & Preservation and reconstruction of spatially localized, boundary-relevant, small-structure, and high-frequency information. & HD95/NSD or other surface measures; thin-structure or small-organ performance; topology/boundary error; controlled high-frequency or resolution perturbations. & Overall DSC alone; visual sharpness alone; a particular convolutional operator. \\
Global contextual reasoning (GC) & Use of long-range, inter-slice, relational, or anatomical-layout information that cannot be resolved from a local neighborhood alone. & Distant/symmetric-structure confusion; long-span relational probes; inter-slice consistency; controlled context ablations; performance stratified by required interaction span. & A large receptive field alone; DSC alone; use of attention or an SSM block. \\
Computational efficiency (CE) & Realized resource use for a specified implementation and input geometry. & Latency, peak memory, energy, throughput, FLOPs, and parameter count measured under a fixed model scope, precision, hardware/counter, and preprocessing policy. & Asymptotic complexity alone; parameter count alone; FLOPs measured under unmatched inputs. \\
Cross-domain generalization (CG) & Preservation of task performance under prospectively held-out domains that share the clinical task, label ontology, and metric-aggregation rule. & Held-out-site/modality performance, degradation relative to in-domain performance, calibration under shift, and prespecified external-domain tests. & Number of datasets evaluated; transfer across incompatible label spaces; zero-shot breadth without a matched task. \\
\bottomrule
\end{tabularx}
\end{table*}

Figures~\ref{fig:taxonomy} and \ref{fig:redistribution_map} should be read in sequence. Figure~\ref{fig:taxonomy} answers \emph{what mechanism defines each family}; Figure~\ref{fig:redistribution_map} answers \emph{which limitation is targeted and which further constraint should be tested}. Method counts, benchmark values, and failure diagnostics are intentionally kept out of these overview figures and reported in the corresponding tables.
\begin{table*}[t]
\centering
\caption{Retrospective coverage and evidence inventory for the structured narrative synthesis. The table describes the included material and its documentation; it does not reconstruct a database-search or screening protocol.}
\label{tab:protocol}
\footnotesize\renewcommand{\arraystretch}{1.16}\setlength{\tabcolsep}{5pt}
\begin{tabularx}{\textwidth}{>{\raggedright\arraybackslash}p{3.1cm}>{\raggedright\arraybackslash}X}
\toprule
\textbf{Aspect} & \textbf{Retrospective description} \\
\midrule
Material represented & Primary medical segmentation papers and author preprints; foundational mechanism papers and prior reviews for context; released implementations where used to clarify a reported metric or mechanism. \\
Mechanism coverage & CNN/U-Net, transformer and CNN--transformer hybrids, SSM/Mamba, SAM/medical foundation models, and adaptive routing/architecture search. \\
Coverage boundary & Retrospective synthesis: work publicly available by 30 June 2026. A targeted contextual update checked on 20 November 2026 is identified separately in Section~\ref{subsec:context_update_2026}; it does not alter the retrospective corpus. \\
Counted material & Named medical segmentation method families discussed in the manuscript. Configurations are grouped under the parent method; a method discussed in several sections is counted once. \\
Contextual material & General-purpose mechanisms, datasets, metric implementations, and prior reviews remain contextual sources and are excluded from the medical-method total. \\
Inventory fields & S1 records method identity, family, grouped variants, citation, manuscript location, and inclusion rationale. S2 records numerical sources, inspected versions, checked fields, and unresolved details. \\
Family assignment & Each method receives one primary architectural family; hybrid mechanisms remain explicit in its description. \\
Evidence use & Values retain their source and configuration. Comparability tiers apply to specified comparisons; unresolved fields do not support pooling, cross-paper normalization, or a causal ranking. \\
\bottomrule
\end{tabularx}
\end{table*}
\subsection{Structured Narrative Synthesis and Corpus Construction}
\label{subsec:sns}

The Structured Narrative Synthesis presented here covers architectural mechanisms and the circumstances under which they are useful. The coverage is purposive, i.e. it aims to represent the different information flow mechanisms rather than providing an estimate of a pooled effect or giving a complete overview of the field. Table~\ref{tab:protocol} gives an overview of the sources and method families covered in the manuscript, as well as the coverage boundary and evidence-recording policy. The accompanying inventories contain the results of the synthesis; they do not reproduce the database queries, the retrieval histories or the screening decisions.

Additional CVPR 2026 and MICCAI 2026 studies are discussed as a separately identified contextual update in Section~\ref{subsec:context_update_2026}. They are excluded from the frozen 95-family method inventory and its configuration-level numerical audit. This targeted update does not extend the retrospective coverage boundary or establish a new systematic screening process.

\noindent\textbf{Method inventory.}
Supplementary Data S1 contains the medical segmentation method families counted in the retrospective synthesis, with a stable record identifier, primary family assignment, citation key, manuscript location, and inclusion rationale. Table~\ref{tab:corpus_audit} is calculated from this inventory. Configurations of a method are recorded as variants and do not silently increase the family count; a method discussed in several sections is counted once. The inventory documents the content of this retrospective synthesis and should not be interpreted as a reconstructed database-search or screening log.

General purpose mechanism papers, datasets, metric implementations and prior surveys are not counted towards the number of medical methods as such. They are, however, cited as context to the medical adaptations discussed. A medical method with several mechanisms is assigned to a single primary family, while the hybrid description is retained in the record. Frequency-aware medical SSM methods discussed in the roadmap are included in the inventory as they provide relevant architectural evidence within the coverage window.
\begin{table}[t]
\centering
\caption{Inventory of named medical segmentation method families. Counts are computed from Supplementary Data S1; configurations are grouped under their parent method and each family is counted once. These counts describe the manuscript's purposive coverage, not field-wide prevalence.}
\label{tab:corpus_audit}
\small\renewcommand{\arraystretch}{1.12}
\begin{tabularx}{\columnwidth}{Xr}
\toprule
\textbf{Primary architectural mechanism} & \textbf{Methods} \\
\midrule
CNN / U-Net & \CorpusCNN \\
Transformer / CNN--transformer hybrid & \CorpusTransformer \\
SSM / Mamba & \CorpusSSM \\
SAM / medical foundation model & \CorpusSAM \\
Adaptive routing / architecture search & \CorpusRouting \\
\midrule
\textbf{Total} & \textbf{\CorpusTotal} \\
\bottomrule
\end{tabularx}
\end{table}

\paragraph{Coverage and documentation limits}
The review is limited by the absence of a contemporaneous database search and screening log for the current review. It is not possible to report on the number of database retrievals, the number of deduplicated records, or the agreement between independent screenings. No formal risk-of-bias ratings are included. The retrospective inventory of included material is given in the manuscript, but this inventory does not establish that all possible material was retrieved, nor that selection was independently duplicated. The number of families included in the manuscript is given, but this does not establish that the material included is scientifically mature, that the evidence is strong, or that the material is representative of the field.

\paragraph{Evidence records and verification scope}
Supplementary Data S2 contains the source checks for the corrected numerical values in the manuscript. It contains the method configuration, task, metric label, source location, and any open accounting details. Verification is reported for the particular value and configuration checked. This does not mean that every experiment in a paper was reproduced, nor does it mean that every reference was fully methodologically assessed. Unknown values are not inferred from familiar benchmarks. Bibliographic transcription and numerical verification are tracked separately.

For any numerical comparison, the evidence record should identify the cohort, patient partition, label ontology, dimensionality, preprocessing and input geometry, initialization and training conditions, metric implementation and aggregation, prompting regime, and computational-accounting scope. When these details are incomplete, a reported value can describe an operating point but cannot by itself establish a pooled effect, normalized ranking, or causal attribution to architecture alone.
\begin{figure*}[t]
\centering
\begingroup
\definecolor{ReviewCNN}{HTML}{0646B5}
\definecolor{ReviewCNNFill}{HTML}{E8F0FF}
\definecolor{ReviewAttention}{HTML}{6B55C9}
\definecolor{ReviewAttentionFill}{HTML}{E5DEFD}
\definecolor{ReviewSSM}{HTML}{4A1FA0}
\definecolor{ReviewSSMFill}{HTML}{EEE3FF}
\definecolor{ReviewSAM}{HTML}{1F5FBF}
\definecolor{ReviewSAMFill}{HTML}{EEF5FF}
\definecolor{ReviewFusion}{HTML}{2E7D32}
\definecolor{ReviewFusionFill}{HTML}{D9F0D5}
\definecolor{ReviewConcern}{HTML}{C00000}
\definecolor{ReviewConcernFill}{HTML}{FFF0F0}
\definecolor{ReviewLine}{HTML}{555555}
\begin{tikzpicture}[
    every node/.style={font=\sffamily\fontsize{9}{10.8}\selectfont,text=black},
    arr/.style={-{Latex[length=2mm]},line width=0.7pt,draw=ReviewLine},
    stage/.style={draw=ReviewCNN,fill=ReviewCNNFill,rounded corners=3pt,
        line width=0.7pt,text width=28.2mm,minimum height=20.2mm,align=center,inner sep=1.5mm},
    decision/.style={stage,draw=ReviewFusion,fill=ReviewFusionFill},node distance=4mm]
\node[stage] (r) {Reported result\\\textbf{source +\\provenance}};
\node[stage,right=of r] (task) {Cohort, task\\and label ontology};
\node[stage,right=of task] (input) {Dimensionality, split\\and input geometry};
\node[stage,right=of input] (metric) {Metric semantics,\\prompts and\\cost scope};
\node[decision,right=of metric] (claim) {Permitted inference\\\textbf{Tier A / B / C}};
\draw[arr] (r)--(task);\draw[arr] (task)--(input);
\draw[arr] (input)--(metric);\draw[arr] (metric)--(claim);
\end{tikzpicture}
\endgroup
\caption{Evidence-admissibility workflow. The workflow specifies the source and protocol checks required for a comparison, including differences in cohort, label space, split, geometry, metric definition, prompting protocol, and computational scope. Tier A permits protocol-local statements for the compatible metric; Tier B supports contextual interpretation; Tier C provides architectural context only. The workflow does not make unmatched studies experimentally equivalent. Drawing code and layout were prepared with assistance from ChatGPT (OpenAI; historical model version not recorded).}
\label{fig:admissibility_flow}
\end{figure*}

\subsubsection{Evidence Provenance and Protocol-Comparability Audit}
\label{subsec:evidence_audit}

ChatGPT (OpenAI) assisted with literature organization and drafting of the source-verification records. The inspected sources, transcribed values, and unresolved fields are documented in S2; the scope of AI assistance is stated in the manuscript-preparation declaration.

Protocol aliasing occurs when results bearing the same dataset name are obtained with different label sets, patient partitions, input dimensionalities, geometries, prompting strategies, metric definitions, or computational-accounting conventions. Figure~\ref{fig:admissibility_flow} summarizes the resulting evidence checks. The inference permitted by a reported result is limited to the protocol fields established in its source record; unresolved fields do not support claims of a strictly matched comparison. For a reported result $r_i$, a comparison record should contain the following tuple:
\begin{equation}
q_i=(\mathcal{D}_i,\mathcal{T}_i,\mathcal{S}_i,\mathcal{G}_i,\mathcal{M}_i,\mathcal{P}_i,\mathcal{C}_i),
\label{eq:evidence_tuple}
\end{equation}
where $\mathcal{D}$ denotes dataset/cohort, $\mathcal{T}$ the task and label ontology, $\mathcal{S}$ the split, $\mathcal{G}$ the dimensionality and input geometry, $\mathcal{M}$ the metric and its reported semantics, $\mathcal{P}$ the prompting or inference regime when relevant, and $\mathcal{C}$ the computational-accounting scope. Missing fields are not imputed from convention.

Table~\ref{tab:comparability_tiers} defines the evidence-use hierarchy. The tiers govern \emph{what claim a row is allowed to support}; they are not scores of model quality. A Tier-A designation applies to a specific comparison and requires compatible patient partitions, task and labels, input and evaluation geometry, metric implementation and aggregation, and relevant inference conditions. Architecture-level causal claims additionally require controlled training and initialization. Tier-B rows remain useful for descriptive context but are excluded from strict numerical ordering whenever a missing or changed field could alter the interpretation. Tier-C rows are retained only as architectural or operating-context evidence.

\begin{table*}[t]
\centering
\caption{Protocol-comparability hierarchy used for quantitative evidence. The hierarchy controls the strength of inference rather than ranking study quality.}
\label{tab:comparability_tiers}
\footnotesize
\renewcommand{\arraystretch}{1.08}

\begin{tabularx}{\textwidth}{>{\raggedright\arraybackslash}p{1.6cm}>{\raggedright\arraybackslash}p{7.1cm}>{\raggedright\arraybackslash}X}
\toprule
\textbf{Tier} & \textbf{Operational criterion} & \textbf{Permitted use in this review} \\
\midrule
A: strict & Same verified patient partition, task and labels; compatible dimensionality, input geometry, metric implementation/units, aggregation, and inference conditions. Within-study ablations must document the held-fixed conditions. & Comparison only for the verified endpoint and protocol. Causal attribution additionally needs matched initialization/training and the intended intervention. Costs need matched scope, inputs, precision, and accounting. \\
B: partial & Core dataset/task is shared, but at least one comparison-relevant field is different, incompletely reported, or not source-verified (e.g., split, input geometry, metric semantics, prompting regime). & Source-specific descriptive context; no strict metric ordering on unresolved fields and no normalized cross-paper score. \\
C: contextual & Different dimensionality, label ontology, prompting/inference regime, or otherwise non-matched evaluation geometry. & Architectural context only; never pooled or normalized with Tier-A/B rows. \\
\bottomrule
\end{tabularx}
\end{table*}

Provenance and comparability are two different things. An original method evaluation and an independent reimplementation of the same method are both primary experimental records of the measurement(s) reported by the authors of these records. A later table that only lists a number of measurements, and for each measurement lists the source of the measurement, is a secondary transcription of the measurements in the table. The unresolved verification marker in the numerical tables does not certify the value of the measurement.

The original paper’s metric semantics are used unless the released evaluation code provides a more explicit definition of the metrics. For example, the use of HD95 in a paper’s results may be abbreviated as “HD” but any such abbreviation would have to be explicitly stated. Decoder-only cost would not be substituted for whole-model cost, and measurements on different input tensors or datasets would not be combined into a single operating point. Missing variance would not be reconstructed from improvements over other baselines or even from decimal precision.

This is a set of rules to keep the comparison of results within reasonable bounds, while not being so restrictive as to hide missing protocol details behind false numerical precision. Results that cannot be strictly compared are reported as operating context.




\section{Datasets and Evaluation Metrics}

Interpretation of segmentation results depends not only on the models used but also on the benchmark and evaluation protocol. Datasets for medical imaging contain different anatomy, imaging modalities, dimensions, class distributions, image qualities and even the purposes for which they were acquired. Thus, different datasets can reveal different weaknesses of a model. Large abdominal organs can be segmented well by a model while small lesions or tumor margins with irregular shape and low contrast can be missed. Models trained on images from another institution can also fail on images from a different institution. Therefore, the choice of dataset for evaluation determines which aspects of the segmentation are being evaluated.

Abdominal multi-organ CT is a very common setting for the evaluation of the anatomical consistency of models on very different structures in terms of size, shape and even location. However, many papers use the name of the dataset Synapse as if it were the name of the challenge BTCV and vice versa. Synapse is the name of the repository where the data of the challenge BTCV are stored and distributed. The challenge was released in the MICCAI 2015 Multi-Atlas Labeling Beyond the Cranial Vault challenge~\citep{landman2015miccai}.

Most of the reported results on Synapse and BTCV datasets are based on different evaluation protocols. For the Synapse dataset, the most popular setting is the 8-organ setting with 18 training images and 12 test images used in TransUNet~\citep{chen2021transunet}. Other settings include a 13-organ setting with different splits and even different pre-processing methods. Unless specified otherwise, the Synapse dataset in this review refers to the 8-organ setting used in TransUNet~\citep{chen2021transunet}.

The large, high-contrast organs in the Synapse dataset can be segmented well by most models. However, success for the smaller, less distinct organs in Synapse depends on a model’s ability to correctly identify the boundaries of the organ as well as to place the organ in the context of the surrounding anatomy. TransUNet~\citep{chen2021transunet}, Swin-Unet~\citep{cao2022swinunet}, and UNETR~\citep{hatamizadeh2022unetr} are three examples of models that have been reported to achieve good results for Synapse. The reported scores for these models are listed in Table~\ref{tab:transformer_medseg} along with information about the dimensionality of the images used, the split of the data into training and test sets, and details of the implementation of the model.

A central challenge in BraTS~\citep{menze2015multimodal} is the variation in tumor appearance, extent, and morphology across patients and scans. Some tumors have well-defined margins, whereas others contain irregular or infiltrative regions, and clinically relevant subregions differ substantially in size and contrast. Whole-tumor segmentation may benefit from broad spatial context, while delineating the enhancing tumor or tumor core also requires sensitivity to small, irregular regions. These distinctions should inform the interpretation of aggregate segmentation scores.

In ACDC~\citep{bernard2018deep}, the geometry of the left ventricle, right ventricle, and myocardium varies between end-diastole and end-systole. Evaluations should therefore account for both anatomy and cardiac phase. Strong Dice scores reported for transformer- and state-space-based models do not, on their own, isolate the effect of contextual modeling: temporal sampling, preprocessing, loss design, and training strategy also influence the comparison.

Boundary-sensitive two-dimensional datasets such as the skin lesion benchmarks (e.g. ISIC~\citep{codella2019skin}) contain irregular contours, weak contrast, color variation and even artifacts. Similarly, the polyp segmentation benchmarks (e.g. Kvasir-SEG~\citep{jha2019kvasir} and CVC-ClinicDB~\citep{bernal2015wm}) contain images with varying scale, illumination, viewpoint, specular reflection and motion blur. In such datasets, accurate localization of the object is not sufficient; the model must also be able to accurately delineate the boundary of the object. Models that retain high-resolution features and use multi-scale decoding as well as models with explicit boundary supervision have reported Dice scores of over 90\% in such datasets, emphasizing the importance of local as well as high-frequency features in such models.

The overlap scores reported for state-space architectures (especially Mamba-based ones) on Synapse, ACDC and related datasets are similar to those of transformer-based methods. While the state-space recurrence of such models scales linearly with the sequence length (instead of the quadratic token-interaction cost of full self-attention), the practical runtime and memory costs depend on several factors, including the image dimensionality, the scan direction, the number of scan paths, the kernel implementation, the parallelization and the hardware. In order to preserve spatial details that may be lost when the image or volume is serialized, many medical SSM architectures include convolutional branches or local operators. 
Most medical segmentation studies report a small set of overlap and boundary metrics. The Dice Similarity Coefficient (DSC) remains the most widely used:

\[
\mathrm{DSC}
=
\frac{2|P \cap G|}{|P| + |G|},
\]

where $P$ denotes the predicted segmentation and $G$ the reference annotation. Dice summarizes the agreement between the predicted and reference regions and is less dominated by background pixels than conventional accuracy. It is therefore well suited to imbalanced segmentation tasks. However, two predictions with similar Dice scores may differ considerably in the location and clinical importance of their errors.

Intersection over Union (IoU) measures the fraction of the combined predicted and reference regions that is shared:

\[
\mathrm{IoU}
=
\frac{|P \cap G|}{|P \cup G|}.
\]

Dice and IoU are mathematically related and usually rank models similarly, although IoU is numerically more stringent for the same prediction. Both are region-based measures and provide limited information about the spatial distribution of boundary errors.

Boundary disagreement is commonly measured using the 95th-percentile Hausdorff distance (HD95). Let the directed distances between boundary point sets $X$ and $Y$ be

\[
\begin{aligned}
\Delta(X,Y)&=\{\min_{y\in Y}d(x,y):x\in X\},\\
\Delta(Y,X)&=\{\min_{x\in X}d(y,x):y\in Y\}.
\end{aligned}
\]

Two conventions called HD95 occur in evaluation software. One takes the larger directed percentile,
\[
\begin{aligned}
H_{95}^{\mathrm{dir}}(X,Y)=\max\{&Q_{0.95}(\Delta(X,Y)),\\
&Q_{0.95}(\Delta(Y,X))\}.
\end{aligned}
\]
The MedPy implementation instead takes a percentile after concatenating the two sets of distances~\citep{medpy_hd95_verified},
\[
H_{95}^{\mathrm{pool}}(X,Y)=Q_{0.95}\bigl(\Delta(X,Y)\mathbin{\|}\Delta(Y,X)\bigr),
\]
where $\|$ denotes concatenation, retaining sample multiplicities. These values need not agree. The released TransUNet and MISSFormer evaluation utilities call MedPy's \texttt{hd95}~\citep{transunet_code_verified,missformer_code_verified}; a paper's short label alone is therefore insufficient to identify the implementation. Surface extraction, spacing, empty-mask handling, and case/class aggregation must also be documented. In particular, a distance is in physical units only when the corresponding spacing is applied. Percentile boundary measures reduce sensitivity to isolated extreme errors relative to the maximum Hausdorff distance, but do not remove the need to inspect clinically important local failures.

The distinction between overlap and boundary errors is of clinical relevance as high Dice scores can be achieved while local contour errors remain undetected near tumor margins, vessels, or treatment borders. This kind of error only minimally changes the total overlap, yet can affect radiotherapy planning, surgical guidance, and even volumetric measurements. Conversely, small structures can obtain low Dice scores after minor boundary shifts, despite the majority of their contour still being anatomically plausible. Thus, overlap and boundary metrics should be analyzed jointly.

Metric definitions alone do not guarantee comparability. Depending on whether results are average by case or by class, how empty masks are handled, whether background is included, and whether surface distances are measured in pixels or physical units, reported values can differ. Furthermore, the terms Hausdorff distance, HD95, mean surface distance, and normalized surface Dice are often used inconsistently. Therefore, whenever comparing numerical results, these details should be reported explicitly.

Cross-paper comparisons are further complicated by differences in data partitioning, resampling, crop size, input dimensionality, normalization, augmentation, pretraining, loss functions, post-processing, and evaluation code. Results obtained computationally are also highly sensitive to model configuration, batch size, numerical precision, hardware, and operator-counting conventions. Thus, a higher benchmark score may indicate a superior architecture, a more favorable experimental setting, or both.

Therefore, different datasets should be considered as complementary stress tests rather than interchangeable leaderboards. Abdominal CT evaluates multi-organ consistency among structures of varying scales. BraTS evaluates heterogeneous tumor compositions and recovery of small subregions. ACDC introduces phase-dependent cardiac anatomy; ISIC and polyp datasets emphasize irregular and weakly defined boundaries. No single benchmark covers all of these requirements.

A convincing evaluation should therefore utilize both region-overlap and boundary measures; clearly document the underlying protocol; and test in multiple anatomical or imaging settings. External domain testing is particularly critical for claims concerning robustness or generalization. Architectural conclusions derived from a single dataset, metric, or data split should be interpreted as evidence from that specific setting rather than proof of general superiority.
\begin{figure*}[!t]
\centering
\includegraphics[width=\textwidth]{assets/architecture_mechanisms.pdf}
\caption{Six information-flow mechanisms used in medical image segmentation. (a) U-Net connects encoder and decoder features at matched spatial scales~\citep{ronneberger2015unet}. (b) UNet++ adds nested refinement along skip pathways~\citep{zhou2018unetplusplus}. (c) The inter-stage mechanism of DoubleU-Net conditions the second network on the product of the image and the first predicted mask; its final mask combination and internal layers are omitted~\citep{jha2020doubleu}. (d) TransUNet combines convolutional features, token self-attention, and decoder skip pathways~\citep{chen2021transunet}. (e) Selective state-space processing propagates a recurrent state with input-dependent updates and readouts~\citep{gu2023mamba}. (f) Learned input-dependent weights combine expert outputs, as exemplified by skip-pathway mixtures in TA-MoSC~\citep{luo2025tamosc_verified}. Expert number, sparsity, stage count, and layer placement are model-specific. These original schematics isolate mechanisms; arrows indicate data or mixture-weight flow, not a performance ordering. Drawing code and layout were prepared with assistance from ChatGPT (OpenAI; historical model version not recorded).}
\label{fig:uneteval}
\end{figure*}

% BEGIN AUTHOR FIGURE fig:original_architecture_panels
Figure~\ref{fig:original_architecture_panels} expands these mechanisms into representative architectural layouts, with the routing panel showing the generic sparse-expert design pattern.

\begin{figure*}[!t]
\centering
\setlength{\tabcolsep}{0pt}
\newcommand{\OriginalArchitecturePanel}[3]{%
  \begin{minipage}[t]{0.485\textwidth}
    \centering
    \begin{minipage}[c][#3][c]{\linewidth}
      \centering
      \includegraphics[width=\linewidth,height=#3,keepaspectratio]{#1}
    \end{minipage}\par\vspace{1.2mm}
    {\footnotesize\bfseries #2\par}
  \end{minipage}%
}
\OriginalArchitecturePanel{figures/Unet arch.png}{(a) U-Net}{0.240\textheight}
\hfill
\OriginalArchitecturePanel{figures/UNet++.png}{(b) UNet++}{0.240\textheight}
\par\vspace{2.0mm}
\OriginalArchitecturePanel{figures/DU_corrected.pdf}{(c) DoubleU-Net}{0.205\textheight}
\hfill
\OriginalArchitecturePanel{figures/trans.drawio.png}{(d) TransUNet}{0.205\textheight}
\par\vspace{2.0mm}
\OriginalArchitecturePanel{figures/SSMMMA.png}{(e) SSM / Mamba}{0.205\textheight}
\hfill
\OriginalArchitecturePanel{figures/moef.png.png}{(f) Sparse MoE layer (generic)}{0.205\textheight}
\par\vspace{1.0mm}
\caption{Detailed architecture panels complementing the mechanisms in Figure~\ref{fig:uneteval}. (a) U-Net~\citep{ronneberger2015unet}; (b) UNet++~\citep{zhou2018unetplusplus}; (c) DoubleU-Net~\citep{jha2020doubleu}; (d) TransUNet~\citep{chen2021transunet}; (e) selective state-space processing~\citep{gu2023mamba}; and (f) a generic sparsely routed mixture-of-experts layer~\citep{fedus2022switch,zoph2022st}. The panels show encoder--decoder pathways, nested feature reuse, cascaded mask conditioning, token attention, recurrent state updates, and expert selection with weighted aggregation, respectively. Panel (f) illustrates the sparse-MoE design pattern; the task-adaptive skip mixture discussed in Section~\ref{sec:cnn} has its own routing implementation. These diagrams describe information flow and do not constitute a matched performance comparison. Panel (c) was revised using vector drawing code prepared with assistance from ChatGPT (OpenAI; historical model version not recorded).}
\label{fig:original_architecture_panels}
\end{figure*}
% END AUTHOR FIGURE fig:original_architecture_panels


\section{The Convolutional Blueprint: Representation Formation and Reconstruction}
\label{sec:cnn}


Figure~\ref{fig:uneteval} places U-Net at the starting point of the architectural progression examined in this review. Its significance extends beyond its historical role. U-Net established the encoder--decoder information pathway that continues to shape later CNN, transformer, SSM, and foundation-model architectures. The encoder progressively reduces spatial resolution while building more abstract representations, whereas the decoder restores a dense prediction. Skip connections create a parallel route through which high-resolution encoder features can bypass the bottleneck and recover spatial detail lost during downsampling.

For encoder feature $E_l$ and decoder feature $D_l$ at level $l$, the canonical U-Net update can be written as

\begin{equation}
\label{eq:skip}
\begin{aligned}
E_l &= \sigma(W_l * E_{l-1}),\\
D_l &= \sigma\!\left(W'_l * \mathrm{Cat}\big(\mathrm{Up}(D_{l+1}),E_l\big)\right).
\end{aligned}
\end{equation}

Equation~\ref{eq:skip} exposes two separate design problems. \emph{Representation formation} concerns what the encoder can learn: convolution captures local texture, edges, and shape cues well, while distant anatomical relationships emerge only gradually through depth, pooling, and receptive-field growth. \emph{Representation transfer} concerns what reaches the decoder. Direct concatenation provides no explicit semantic-alignment mechanism. Subsequent decoder operations must learn how to combine shallow appearance cues with deeper task-level features. Many U-Net variants can be read as modifications to one of these two processes rather than replacements of the basic encoder--decoder design.

\subsection{Representation Formation in the U-Net Lineage}

The development of convolutional U-shaped architectures follows a clear pattern. Residual and recurrent variants improve optimization and repeated local refinement. Short residual shortcuts have been studied directly in biomedical segmentation alongside long encoder--decoder skips~\citep{drozdzal2016skip_verified}. The Residual U-Net results below refer specifically to the medical reimplementation evaluated by \citet{azad2024unet_verified}, rather than to experiments on road extraction. Dense architectures encourage feature reuse and strengthen gradient flow. UNet++~\citep{zhou2018unetplusplus} inserts intermediate transformations along the skip pathways to reduce the semantic gap between encoder and decoder features. Attention U-Net~\citep{oktay2018attentionunet} changes skip fusion from unconditional transfer to task-dependent selection. Multi-scale designs, including U-Net 3+~\citep{review95_huang2020unet3plus} and Double U-Net~\citep{jha2020doubleu}, expose the decoder to features from a wider range of spatial resolutions and support coarse-to-fine reconstruction. Three-dimensional variants extend the same design principles to volumetric data, while nnU-Net~\citep{isensee2021nnu} shows that strong performance depends not only on network structure, but also on preprocessing, patch sampling, augmentation, loss design, and training configuration.

PMRNet~\citep{kang2026pmrnet_cvpr} provides a further lightweight design, combining physics-informed multi-scale convolution, diffusion-based contextual propagation, and boundary-aware decoding. Its reported computational and segmentation outcomes remain specific to the authors' evaluation settings and are not pooled with the heterogeneous records in Table~\ref{tab:unet_compare}.

These developments mark the main changes within the convolutional lineage. CNN variants improve how local evidence is extracted, preserved, and reconstructed, while making clear that feature fusion is a separate architectural problem. As skip pathways become more selective and more heavily transformed, the question shifts from whether encoder features should be transferred to which features should be transferred, from which scale, and under what form of conditioning.

Table~\ref{tab:unet_compare} is therefore kept to representative convolutional and U-shaped architectures. Section~\ref{subsec:skip_evolution} treats skip-connection design as an information-flow problem in its own right. Transformer- and SSM-based methods are considered separately, where their long-range context mechanisms can be discussed without repeating the U-Net taxonomy.

\begin{table*}[!t]
\centering
\fontsize{7.5}{9}\selectfont
\renewcommand{\arraystretch}{1.05}
\setlength{\tabcolsep}{3pt}
\begin{tabularx}{\textwidth}{>{\raggedright\arraybackslash}p{2.2cm}>{\raggedright\arraybackslash}p{2.0cm}>{\raggedright\arraybackslash}p{1.65cm}>{\raggedright\arraybackslash}p{1.3cm}>{\raggedright\arraybackslash}p{1.65cm}>{\raggedright\arraybackslash}p{1.0cm}>{\raggedright\arraybackslash}p{.95cm}>{\raggedright\arraybackslash}X}
\toprule
\textbf{Method} & \textbf{Design change} & \textbf{Dataset} & \textbf{DSC (\%)} & \textbf{Boundary metric} & \textbf{Params (M)} & \textbf{FLOPs (G)} & \textbf{Numerical source / configuration} \\
\midrule

U-Net \citep{ronneberger2015unet}
& Direct encoder--decoder skips
& Synapse CT
& 76.85
& HD 39.70
& 1.95
& 8.2
& Later medical implementation: \citet{azad2024unet_verified}\textsuperscript{\dag}, Tables II(c), III; 2D $224\times224$. \\

UNet++ \citep{zhou2018unetplusplus}
& Nested skip refinement
& Synapse CT
& 76.91
& HD 36.93
& 9.16
& 26.72
& Later medical implementation: \citet{azad2024unet_verified}\textsuperscript{\dag}, Tables II(c), III; 2D $224\times224$. \\

Dense U-Net \citep{cai2020dense}
& Dense feature reuse
& In-vivo multiphoton microscopy
& 90.60
& NR
& NR
& NR
& Original cellular segmentation study; the number does not describe abdominal CT. \\

H-DenseUNet \citep{li2018h}
& Hybrid 2D--3D convolution
& LiTS 2017
& Liver 96.1; lesion 72.2
& NR
& NR
& NR
& Primary Table II, Dice per case. Global Dice is a different aggregation and is not substituted. \\

Residual U-Net (medical reimplementation) \citep{azad2024unet_verified}
& Residual feature propagation
& Synapse CT
& 76.95
& HD 38.44
& 13.04
& 62.06
& \citet{azad2024unet_verified}\textsuperscript{\dag}, Tables II(c), III; 2D $224\times224$. This is the evaluated medical configuration. \\

R2U-Net \citep{alom2018recurrent}
& Recurrent residual refinement
& Synapse CT
& 77.12
& HD 32.12
& NR
& NR
& Later comparison in \citet{zhang2024multi}\textsuperscript{\dag}; the source calls the boundary metric HD. \\

DoubleU-Net \citep{jha2020doubleu}
& Cascaded refinement
& CVC-ClinicDB
& 92.39
& NR
& NR
& NR
& Primary Table III: Dice 0.9239, expressed here as a percentage. Dataset-specific result; no pooled polyp score is implied. \\

KiU-Net \citep{valanarasu2021kiu,review95_valanarasu2020kiunet}
& Undercomplete and overcomplete paths
& Brain ultrasound
& 89.43
& NR
& 0.29
& NR
& Original conference report, Table 1; mean Dice is shown. The reported dispersion is not transferred without its aggregation context. \\

3D U-Net \citep{cciccek20163d}
& Volumetric convolution
& BraTS 2020 validation
& ET 68.76; WT 84.11; TC 79.06
& HD95: ET 50.983; WT 13.366; TC 13.607
& 16.21
& 1670
& Plain 3D U-Net baseline in TransBTS\textsuperscript{\dag} \citep{reviewfix_wang2021transbts}, Table 2 and Section 3.2; cost accounting is source-specific. Residual 3D U-Net is a separate row in that source. \\

nnU-Net \citep{isensee2021nnu}
& Self-configuring pipeline
& Synapse, 8 organs
& 86.21
& HD95 10.91
& 30.64
& 410.11
& Later rebenchmark: \citet{lumetti2025taming}\textsuperscript{\dag}, Tables 3 and 5. Mean results; costs apply to the Synapse implementation. \\

CBAM-U-Net \citep{zulaikha2023liver}
& Channel and spatial attention
& Liver CT
& NV
& NV
& NV
& NV
& Feature recalibration is retained as the architectural contribution; no unverified configuration-specific measurement is displayed. \\

\bottomrule
\end{tabularx}
\caption{Representative convolutional segmentation designs with numerical sources separated from architectural origins. Scores retain the source metric and reported precision. The 3D U-Net row uses the plain baseline, not the residual baseline. Different datasets, label spaces, aggregations, and implementations prevent a common ranking. NR, NV, and the provenance marker follow the convention defined in Table~\ref{tab:transformer_medseg}.}
\label{tab:unet_compare}
\end{table*}


\subsection{Skip-Connection Evolution: From Transfer to Representation Alignment}
\label{subsec:skip_evolution}

Fixing the encoder makes representation transfer an additional design variable. Direct concatenation supplies high-resolution evidence and leaves compatibility to the decoder. The skip literature introduces more explicit mechanisms: UNet++ inserts alignment transformations, Attention U-Net introduces task-conditioned gating, UCTransNet and UDTransNet model cross-scale/channel compatibility, and TA-MoSC treats fusion as conditional routing. The sequence is therefore \emph{direct transfer $\rightarrow$ semantic alignment $\rightarrow$ selective gating $\rightarrow$ cross-scale interaction $\rightarrow$ conditional allocation}. Table~\ref{table:skip_connections} isolates this interface-level progression from changes to the backbone itself.

Skip redesign can improve the exchange of retained features and can form new relationships through attention or other decoder operations. It cannot recover input evidence that has already been discarded. This motivates examining representation formation and transfer together when testing long-range context, spatial fidelity, and deployment cost.
\begin{table*}[t]
\centering\footnotesize
\setlength{\tabcolsep}{4pt}\renewcommand{\arraystretch}{1.14}
\begin{tabularx}{\textwidth}{>{\raggedright\arraybackslash}p{2.65cm}>{\raggedright\arraybackslash}p{2.5cm}>{\raggedright\arraybackslash}p{3.4cm}>{\raggedright\arraybackslash}p{3.5cm}>{\raggedright\arraybackslash}X}
\toprule
\textbf{Method} & \textbf{Fusion mechanism} & \textbf{Conditioning signal} & \textbf{Intended role} & \textbf{Constraint to test} \\
\midrule
U-Net \citep{ronneberger2015unet} & Direct concatenation & No explicit selection gate; learned decoder operations & Preserve spatial features for reconstruction & Compatibility of encoder and decoder features \\
UNet++ \citep{zhou2018unetplusplus} & Dense nested skip pathways & Learned transformations of intermediate features & Refine features before transfer & Activation memory and redundant paths \\
Attention U-Net \citep{oktay2018attentionunet} & Attention gate & Current encoder features and a decoder gating signal & Weight spatially relevant features & Gate errors and added execution cost \\
UCTransNet \citep{wang2022uctransnet} & Channel attention & Current features across encoder scales & Align and combine cross-scale information & Attention cost and transfer quality \\
UDTransNet \citep{wang2024udtransnet} & Spatial--channel attention & Current encoder and decoder representations & Improve compatibility at the fusion interface & Attention cost and sensitivity to feature mismatch \\
TA-MoSC \citep{luo2025tamosc_verified} & Expert routing & Input-pooled features with stage-specific routing & Select among skip-path transformations & Routing calibration, utilization, and execution overhead \\
\bottomrule
\end{tabularx}
\caption{Mechanisms used for skip-path transfer. Conditioning is described directly instead of assigning an ordinal adaptivity score. Intended roles and candidate constraints are hypotheses to test; this table does not establish a performance or cost ordering.}
\label{table:skip_connections}
\end{table*}

PSPNet-based prostate MRI segmentation provides a further task-specific convolutional design in the inventory~\citep{yan2021psp}; its results are not treated as a matched comparison with the abdominal U-Net records.

\section{Architecture-Centric Analysis of Transformer-Based Medical Image Segmentation Methods}


The preceding section showed that increasingly sophisticated fusion can improve the exchange of encoder and decoder information, but cannot recover input evidence discarded by the encoder, although attention or other decoder operations can form new relationships among retained features. Transformers address this upstream limitation by replacing exclusively local aggregation with content-dependent interaction. For queries $Q$, keys $K$, and values $V$, self-attention is
\begin{equation}
\mathrm{Attention}(Q,K,V) =
\mathrm{softmax}\left(\frac{QK^\top}{\sqrt{d_k}}\right)V.
\end{equation}
In segmentation, this mechanism allows distant regions to exchange information directly and is therefore well suited to tasks in which anatomical identity depends on spatial configuration rather than local appearance alone.

The transformer literature in medical segmentation is best organized around three design questions. First, \emph{where should global interaction be inserted?} Hybrid systems retain convolutional stages and place attention at the bottleneck or in parallel branches, whereas pure transformer encoders replace convolutional representation learning more extensively. Second, \emph{how should interaction scale to high-resolution and 3D inputs?} Hierarchical, windowed, axial, deformable, and sparse attention trade complete pairwise interaction for tractable computation. Third, \emph{how should global representations be returned to dense spatial prediction?} Multi-scale decoders, convolutional reconstruction stages, boundary-aware modules, and cross-branch fusion address the same alignment problem identified in the skip-connection section, but now between representations formed under different inductive regimes.

The transformer comparison is therefore organized around how each design balances global interaction, spatial precision, and computational tractability. Tables~\ref{tab:transformer_medseg} and~\ref{tab:transformer_medseg_cont} give the method-level evidence; the subsections move from hybrid designs to hierarchical attention, multi-scale fusion, and more efficient interaction.


The progression begins with hybrid models that add global interaction to a convolutional scaffold, then moves toward transformer-dominant encoders, hierarchical attention, and specialized mechanisms for fusion and efficiency.


\subsection{From Hybrid CNN--Transformer Models to Pure Transformer Architectures}

Early transformer-based segmentation models were hybrids that kept a convolutional stem to supply the locality bias that pure attention lacks. In TransUNet~\citep{chen2021transunet}, a CNN extracts low-level features and a Vision Transformer (ViT)~\citep{dosovitskiy2021vit} encoder models long-range dependencies among the resulting tokens; TransBTS~\citep{reviewfix_wang2021transbts} places a transformer bottleneck inside a 3D CNN so that spatially separated tumour subregions can interact directly. UNETR~\citep{hatamizadeh2022unetr} removes the CNN encoder altogether, applying a ViT to 3D patches and drawing multi-scale skip connections from intermediate transformer layers. The progression trades convolutional locality for a more uniform global representation; the changed spatial inductive bias motivates testing the roles of decoder design, initialization, and data scale to recover the local structure that convolution supplied by construction.

\subsection{Hierarchical and Efficient Transformer Architectures}

Hierarchical designs restrict the interaction graph so that attention remains tractable at high resolution. Swin UNETR~\citep{hatamizadeh2022swinunetr} uses shifted-window attention, so that attention cost grows linearly with total token count for a fixed window volume and context crosses windows only through successive shifted layers. nnFormer~\citep{zhou2023nnformer} interleaves local and global volume attention within 3D blocks, and CoTr~\citep{xie2021cotr} uses deformable attention to attend to a sparse set of sampled locations rather than the full feature map. These designs control cost by restricting the number or structure of interactions; the remaining question is which long-range relations are lost when that decision does not match the anatomy.

\subsection{Fusion of Global and Local Representations}


The skip-connection section established that fusion quality depends on representation compatibility. In transformer systems this problem becomes more pronounced because convolutional and attention branches encode complementary but statistically different information. Consequently, the central design problem is not simply to concatenate ``local'' and ``global'' features, but to align them across scale, channel, and semantic abstraction.

TransFuse~\citep{reviewfix_zhang2021transfuse} addresses this by maintaining parallel CNN and transformer branches and using dedicated BiFusion modules rather than assuming direct compatibility. Other architectures solve the same problem through hierarchical multi-scale representations: MISSFormer~\citep{reviewfix_huang2021missformer} and MERIT~\citep{reviewfix_rahman2024merit} aggregate features across resolutions, while PVT-CASCADE~\citep{reviewfix_rahman2023cascade} uses a pyramid representation to preserve scale diversity during decoding. Boundary-oriented designs such as FAT-Net~\citep{wu2022fatnet} further emphasize that global semantic correctness must be converted back into high-frequency spatial structure.

Thus, transformer fusion should be interpreted as a continuation of the skip-connection problem under a harder representation mismatch. The novelty lies in reconciling features produced by different context-forming mechanisms, rather than merely increasing the number of fusion operations.


\subsection{Lightweight and Efficient Transformer Designs}

Several designs target the quadratic cost of attention more directly. Axial-DeepLab provides general-domain provenance for position-sensitive axial attention, which factorizes 2D attention along height and width~\citep{wang2020axial_verified}. Its panoptic-segmentation experiments are distinct from medical U-Net evaluation. Valanarasu et al. evaluate an axial-attention U-Net baseline and the gated Medical Transformer (MedT) on medical datasets~\citep{valanarasu2021medt_verified}; the MedT row in Table~\ref{tab:transformer_medseg} reports the latter model. MedT combines gated axial attention with a local--global training strategy. LeViT-UNet~\citep{xu2023levitunet} uses a lightweight transformer encoder for speed, UTNet~\citep{gao2021utnet} inserts efficient attention into U-Net stages, DAE-Former~\citep{reviewfix_azad2022daeformer} combines efficient spatial and channel attention, and MDFormer~\citep{li2025mdformer} uses dense dilated attention to enlarge the receptive field at fixed resolution. These designs reduce the computational bottleneck by approximating or restricting interaction, which relocates the question to which dependencies the approximation preserves.
\begin{table*}[!t]
\centering
\fontsize{8}{9.5}\selectfont
\renewcommand{\arraystretch}{1.1}\setlength{\tabcolsep}{3pt}
\begin{tabularx}{\textwidth}{>{\raggedright\arraybackslash}p{2.2cm}>{\raggedright\arraybackslash}p{1.75cm}>{\raggedright\arraybackslash}p{3.1cm}>{\raggedright\arraybackslash}p{.9cm}>{\raggedright\arraybackslash}p{.95cm}>{\raggedright\arraybackslash}X}
\toprule
\textbf{Method / variant} & \textbf{Dataset / task} & \textbf{Reported outcome} & \textbf{Params (M)} & \textbf{FLOPs (G)} & \textbf{Mechanism and source context} \\
\midrule
TransUNet \citep{chen2021transunet} & Synapse multi-organ CT & DSC 77.48\%, HD 31.69 & 105.28$^{\dagger}$ & 80.68$^{\dagger}$ & Hybrid CNN--ViT encoder; CNN extracts low-level spatial features while ViT models long-range dependencies; U-Net-like decoder restores resolution. Cost source: \citep{azad2024unet_verified}. \\

UNETR \citep{hatamizadeh2022unetr}$^{\dagger}$ & ACDC / MSD BrainTumor / Synapse & ACDC DSC 88.72\%; BT DSC 70.92\%, HD95 9.04; Synapse DSC 78.10\%, HD95 19.15 & 92.49 & 75.76 & ViT encoder with intermediate-layer skips to a convolutional decoder. Numerical source: \citet{lumetti2025taming}, Tables 2--5; cost is for Synapse. \\

Swin UNETR \citep{hatamizadeh2022swinunetr}$^{\dagger}$ & ACDC / MSD BrainTumor / Synapse & ACDC DSC 91.36\%; BT DSC 84.07\%, HD95 6.77; Synapse DSC 83.64\%, HD95 11.02 & 62.83 & 384.20 & Hierarchical shifted-window encoder with convolutional reconstruction. Numerical source: \citet{lumetti2025taming}, Tables 2--5; cost is for Synapse. \\

nnFormer \citep{zhou2023nnformer}$^{\dagger}$ & ACDC / MSD BrainTumor / Synapse & ACDC DSC 91.87\%; BT DSC 86.34\%, HD95 4.05; Synapse DSC 86.56\%, HD95 11.14 & 150.50 & 213.41 & Interleaved local and global volumetric attention with convolutional embedding. Numerical source: \citet{lumetti2025taming}, Tables 2--5; cost is for Synapse. \\

CoTr \citep{xie2021cotr}$^{\dagger}$ & ACDC / MSD BrainTumor / Synapse & ACDC DSC 90.90\%; BT DSC 68.21\%, HD95 9.96; Synapse DSC 84.67\%, HD95 9.35 & 50.12 & 369.22 & Convolutional encoder with sparse deformable attention. Numerical source: \citet{lumetti2025taming}, Tables 2--5; cost is for Synapse. \\

Lightweight TransBTS \citep{reviewfix_wang2021transbts} & BraTS 2019 validation & Dice WT 90.36\%, TC 81.76\%, ET 78.94\%; HD95 WT 6.004, TC 6.173, ET 4.552 & 15.14 & 208 & Reduced TransBTS variant with one transformer layer and a halved feed-forward hidden dimension; costs are reported in the original paper, with counting geometry unspecified in the cited passage. \\

TransFuse-S \citep{reviewfix_zhang2021transfuse} & ISIC2017 / Kvasir / CVC-ClinicDB & ISIC Dice 87.2\%, Kvasir Dice 91.8\%, ClinicDB Dice 91.8\% & 26.3 & NR & ResNet-34 and DeiT-Small branches fused through BiFusion modules; values refer to the small variant; its original report does not give FLOPs. \\

MedT \citep{valanarasu2021medt_verified} & Brain US / GlaS / MoNuSeg & F1 88.84 (Brain US), 81.02 (GlaS), 79.55 (MoNuSeg) & 1.4 (Brain US) & NR & Gated axial attention and local--global training. The parameter count is the Brain-US configuration in supplementary Table 2; it is not transferred to other datasets. \\
Swin-Unet \citep{cao2022swinunet} & Synapse multi-organ CT & DSC 79.13\%, HD 21.55 & 27.17$^{\dagger}$ & 6.16$^{\dagger}$ & Pure Swin Transformer U-shaped architecture; patch merging encoder and patch expanding decoder. Costs: \citet{azad2023dlka_nv}, Table 1. \\

MISSFormer \citep{reviewfix_huang2021missformer} & Synapse, ACDC & Synapse DSC 81.96\%, HD 18.20 & 42.46$^{\dagger}$ & 7.25$^{\dagger}$ & Hierarchical encoder-decoder transformer with enhanced skip connections and multi-scale feature aggregation. Cost source: \citep{azad2024unet_verified}. \\

LeViT-UNet-384 \citep{xu2023levitunet} & Synapse & DSC 78.53\%, HD 16.84 & 52.17 & 25.55 & LeViT encoder with U-Net reconstruction; primary Tables 1--2 of the author preprint report these values and call the boundary metric HD. \\

UCTransNet \citep{wang2022uctransnet} & Synapse, GlaS, MoNuSeg & Synapse DSC 78.23\%, HD 26.75 & Synapse 66.43$^{\dagger}$ & Synapse 32.99$^{\dagger}$ & Channel-transformer skip fusion. Accuracy: original Table 2; costs: later evaluation by \citet{azad2024unet_verified}. \\
UTNet \citep{gao2021utnet} & Cardiac MRI; vendor A to A & Average Dice 88.3\% & 9.53 & NR & Efficient-attention U-Net; original Table 1 reports this same-vendor cardiac-MRI configuration. \\
HiFormer-L \citep{heidari2023hiformer} & Synapse & DSC 80.69\%, HD95 19.14 & 29.52 & 11.88$^{\dagger}$ & CNN and Swin features fused across scales. Large variant; source Tables 2/5. Section 4.3 defines its HD label as 95th-percentile Hausdorff distance. FLOPs at $224\times224$: \citet{tang2025simmpnet_nv}, Table 2. \\
\bottomrule
\end{tabularx}
\begin{minipage}{\linewidth}\footnotesize
\textbf{Numerical provenance.} $^{\dagger}$ identifies a value from an explicitly cited later evaluation or reimplementation, rather than the method's original experiment. A marker on a method name applies to that row's numerical entries; a marker on a cell applies only to that cell. Unmarked retained numbers come from the method authors' cited primary report or version. NR means that no value for the displayed configuration was reported in the inspected source version; it does not assert absence from every later report. NV means that configuration-specific verification remains unresolved because source access, model attribution, or units are insufficient. Neither label denotes zero. Numbers are source-reported central estimates unless an uncertainty term and its definition are explicitly stated.
\textbf{Protocols and units.} Params denotes the source-reported parameter count in millions; total versus trainable parameters must be checked before comparison. GFLOPs retains each source's operation-count convention. A common input shape, batch size, or multiply--add convention is not assumed. Lumetti et al.'s costs are Synapse-specific and its Tables 2--4 describe fold-based summaries. HD and HD95 labels are preserved separately; a shared unit or metric implementation is not inferred from similar numbers. The original TransUNet code calls \texttt{hd95} without voxel spacing \citep{transunet_code_verified}, so its paper's HD label does not by itself resolve physical-distance comparability. DAE-Former follows Table 1 of the cited preprint; other tables in that version contain different values.
\end{minipage}
\caption{Representative transformer methods, part 1: source-specific evaluations and configurations. Further methods appear in Table~\ref{tab:transformer_medseg_cont}. The values are descriptive operating points and do not form a shared leaderboard.}
\label{tab:transformer_medseg}
\end{table*}

\begin{table*}[!t]
\centering
\fontsize{8}{9.5}\selectfont
\renewcommand{\arraystretch}{1.1}\setlength{\tabcolsep}{3pt}
\begin{tabularx}{\textwidth}{>{\raggedright\arraybackslash}p{2.2cm}>{\raggedright\arraybackslash}p{1.75cm}>{\raggedright\arraybackslash}p{3.1cm}>{\raggedright\arraybackslash}p{.9cm}>{\raggedright\arraybackslash}p{.95cm}>{\raggedright\arraybackslash}X}
\toprule
\textbf{Method / variant} & \textbf{Dataset / task} & \textbf{Reported outcome} & \textbf{Params (M)} & \textbf{FLOPs (G)} & \textbf{Mechanism and source context} \\
\midrule
DS-TransUNet-L \citep{lin2022dstransunet} & Kvasir & mDice 91.3\% & NR & NR & Dual-branch Swin Transformer encoder with multi-scale feature fusion. Original Table II; 880/120 training/test split, $384\times384$ inputs. \\

TransDeepLab \citep{reviewfix_azad2022transdeeplab} & Synapse & DSC 80.16\%, HD 21.25 & 21.14 & 16.31$^{\dagger}$ & Transformer-based DeepLab-style segmentation; values follow the original paper's Synapse comparison and parameter report. FLOPs: \citet{azad2023dlka_nv}, Table 1. \\

DAE-Former \citep{reviewfix_azad2022daeformer} & Synapse & DSC 82.63\%, HD 16.39 & 48.07$^{\dagger}$ & 27.89$^{\dagger}$ & Efficient spatial/channel attention and cross-attention skip fusion. Accuracy follows original Table 1; paired costs follow \citet{azad2023dlka_nv}, Table 1. \\

Cascaded MERIT \citep{reviewfix_rahman2024merit} & Synapse, ACDC & Synapse DSC 84.90\%, HD95 13.22; ACDC DSC 91.85\% & 147.86 & 33.31 & Two hierarchical MaxViT backbones with CASCADE decoding and MUTATION training; reported cost uses $256\times256$ and $224\times224$ inputs. \\

PVT-CASCADE \citep{reviewfix_rahman2023cascade} & Synapse & DSC 81.06\%, HD95 20.23 & 35.28$^{\dagger}$ & 6.40$^{\dagger}$ & Pyramid Vision Transformer encoder with cascaded attention decoding; reported Synapse metrics use $224\times224$ inputs and are averaged over five runs. Costs: \citet{azad2023dlka_nv}, Table 1. \\

FAT-Net \citep{wu2022fatnet}$^{\dagger}$ & ISIC2018 & Dice 89.08\% & 30 & NV & Feature adaptive transformer with CNN--transformer fusion. Later reimplementation: \citet{lin2023batformer}, Table V; 2076/520 training/test split. Its reported F=23 is labelled GFLOP/s. \\

FCT \citep{tragakis2023fct} & ACDC cardiac MRI & ACDC Dice 93.02\% (FCT-384) & NV & NV & Fully convolutional transformer; combines convolutional inductive bias with transformer-style global modeling. \\

H2Former \citep{he2023h2former} & ISIC2016 & Dice 92.41\% & 33.71 & 33.56 & Hierarchical hybrid CNN--transformer fusion. Original Table I and Section IV-E; $512\times512$ inputs, five-experiment mean Dice. \\

2D BATFormer \citep{lin2023batformer} & ACDC cardiac MRI & Dice 92.84\% & 1.2 & NV & Boundary-aware transformer; original Table IV, 70/10/20 patient split. Its reported F=4.1 is labelled as an operation rate. \\

SwinPA-Net \citep{du2024swinpanet}$^{\dagger}$ & CVC-ClinicDB / Kvasir-SEG & ClinicDB Dice 95.68\%; Kvasir Dice 95.04\% & 102.87 & NR & Dense multiplicative connections and local pyramid attention. Later $256\times256$, 8:1:1 evaluations: \citet{chen2026pcatransunet_nv}, Tables 1--2. \\
MIST \citep{rahman2024mist} & ACDC/ Synapse & ACDC Dice 92.56\%, Synapse Dice 86.92\%, HD95 11.07 & 95.74 & NR & MaxViT and CAM decoder. Outcomes follow original Tables 1--2; parameters refer to the Synapse ablation. The printed Synapse aggregate differs from the mean of its listed organ scores; S2 records this unresolved discrepancy. \\
HCTNet \citep{he2023hctnet}$^{\dagger}$ & BUSI breast ultrasound & Dice 81.24\% & 22.20 & 38.91 & Hybrid CNN--transformer with cross-attention. Later $224\times224$ evaluation: \citet{yang2026mambacnn_nv}, Tables 1/3; exact BUSI-only split is unspecified. \\

MT-UNet \citep{reviewfix_wang2021mtunet} & Synapse, ACDC & Synapse DSC 78.59\%, HD95 26.59, ACDC DSC 90.43\% & Synapse 79.07$^{\dagger}$ & Synapse 44.75$^{\dagger}$ & Mixed Transformer U-Net combines local--global Gaussian-weighted self-attention and external attention within a U-shaped architecture. Synapse costs at $224\times224$: \citet{tang2025simmpnet_nv}, Table 2 (MixTrans). \\

UNETR++ \citep{shaker2024unetrpp} & Synapse / BTCV / ACDC & Synapse DSC 87.22\%, HD95 7.53; BTCV DSC 83.28\%; ACDC DSC 92.83\% & 42.96 & 47.98 & Efficient paired spatial/channel attention. Costs and Synapse metrics follow primary Table 1; BTCV uses 13 organs. Costs are not transferred from the BraTS configuration. \\

MDFormer \citep{li2025mdformer} & Synapse, ACDC & Synapse DSC 86.88\%, HD95 6.74, ACDC DSC 92.7\% & 37.7 & 47.39 & Multi-scale dense dilated transformer; combines dilated attention and dense aggregation to enlarge receptive field efficiently. \\
\bottomrule
\end{tabularx}

\caption{Representative transformer methods, part 2. Protocol and provenance conventions follow Table~\ref{tab:transformer_medseg}. NR identifies an inspected-source nonreport and NV an unresolved value; an architecture is included for its mechanism even when a numerical record is not verified.}
\label{tab:transformer_medseg_cont}
\end{table*}


\subsection{Task-Specific and Boundary-Aware Architectures}

Transformer-based models have also been adapted to specific medical tasks. For example, BATFormer~\citep{lin2023batformer} introduces boundary-aware attention to improve edge delineation, addressing a common weakness of transformer models in low-contrast regions. SwinPA-Net~\citep{du2024swinpanet} combines a Swin Transformer backbone with dense multiplicative connections and local pyramid attention to aggregate multiscale medical-image features. HCTNet~\citep{he2023hctnet} combines CNN and transformer features for breast-ultrasound lesion segmentation, where speckle and weak lesion margins make local evidence unreliable.
The common point is practical: transformer blocks alone do not solve low contrast, noisy boundaries, or anatomical variability. Task-specific designs usually add an explicit mechanism for one of these failure modes.

Table~\ref{tab:transformer_medseg} summarizes representative transformer-based medical image segmentation methods. Unavailable or unverified configuration-specific values are marked NR or NV and excluded from numerical comparison. Image resolution, data split, augmentation, model variant, and post-processing must be checked before any numerical comparison.
\paragraph{Mechanistic takeaway}


The transformer trajectory is now clearer when read together with the U-Net and skip-connection sections. U-Net establishes a two-route encoder--decoder information flow; skip-connection research improves the transfer and alignment route; transformers modify the representation-formation route by introducing content-dependent long-range interaction. The major transformer families then differ primarily in how they make that interaction computationally feasible and how they reconnect global representations to dense spatial reconstruction.

Full attention permits all-pairs interaction but scales poorly with token count. Hierarchical and windowed attention reduce this cost by restricting the interaction graph and propagating context across layers. Hybrid models preserve convolutional spatial priors but introduce cross-representation alignment problems. Boundary-aware and multi-scale decoders recover local precision, but add specialized modules and additional design complexity. These are not four separate criticisms; they are consequences of a single architectural tension between interaction breadth, spatial fidelity, and computational cost.

The unresolved mechanism-level question is whether long-range dependency modeling can be retained without explicit all-pairs interaction or increasingly elaborate approximations to it. State-space models address this question by replacing pairwise attention with recurrent state propagation.


\section{Mamba-Based Medical Image Segmentation: From Efficient State Propagation to Spatial Serialization}
\label{sec:mamba_review}


The transformer section ends with a computational question: can long-range dependency modeling be retained without an explicit all-pairs token-interaction matrix? State-space models (SSMs), including Mamba, offer one way to address this question through input-conditioned state propagation. This changes the mechanism of contextual interaction but leaves an important spatial problem: medical images are two- or three-dimensional fields, whereas state updates operate on ordered sequences. The key issue is therefore not simply whether an SSM is more efficient than attention, but how anatomy is serialized, scanned, and reconstructed, and whether spatial relationships remain accessible to the segmentation model.





\begin{figure*}[!t]
\centering
\includegraphics[width=0.8\textwidth]{assets/mamba_mechanisms.pdf}
\caption{Design choices for adapting selective state-space processing to medical images. (a) Planar scan orders serialize two-dimensional features, followed by state updates, merging, and spatial reconstruction, as represented by VM-UNet and Mamba-UNet~\citep{ruan2024vm,wang2024mamba}. (b) Volume-oriented sequences expose different spatial directions to recurrent processing, exemplified by the tri-oriented design of SegMamba~\citep{xing2024segmamba}. (c) Local convolution and state processing can be combined within an encoder--decoder pipeline, as in U-Mamba~\citep{ma2024umamba}. (d) Frequency-domain features can complement spatial/state features, as investigated by EM-Net, HybridMamba, and FaRMamba~\citep{verified_chang2024emnet,wu2025hybridmamba,verified_rong2025farmamba}. The spectral operator, fusion location, and serial or parallel arrangement differ across the cited methods; panel (d) expresses this design axis rather than one shared implementation. These original conceptual diagrams omit implementation-specific blocks and do not establish comparative performance. Drawing code and layout were prepared with assistance from ChatGPT (OpenAI; historical model version not recorded).}
\label{fig:mamba_architecture_progression}
\end{figure*}

% BEGIN AUTHOR FIGURE fig:original_mamba_panels
Figure~\ref{fig:original_mamba_panels} provides complementary architectural views of the representative state-space segmentation methods discussed below.

\begin{figure*}[!t]
\centering
\setlength{\tabcolsep}{0pt}
\newcommand{\OriginalMambaPanel}[3]{%
  \begin{minipage}[t]{0.485\textwidth}
    \centering
    \begin{minipage}[c][#3][c]{\linewidth}
      \centering
      \includegraphics[width=\linewidth,height=#3,keepaspectratio]{#1}
    \end{minipage}\par\vspace{1.2mm}
    {\footnotesize\bfseries #2\par}
  \end{minipage}%
}
\OriginalMambaPanel{figures/U_Mamdr.png}{(a) U-Mamba}{0.170\textheight}
\hfill
\OriginalMambaPanel{figures/Segmambadrawio.png}{(b) SegMamba}{0.170\textheight}
\par\vspace{2.0mm}
\OriginalMambaPanel{figures/VMUNETdrawio.png}{(c) VM-UNet}{0.170\textheight}
\hfill
\OriginalMambaPanel{figures/MambaUNetdrawio.png}{(d) Mamba-UNet}{0.170\textheight}
\par\vspace{2.0mm}
\OriginalMambaPanel{figures/SwinUMambadrawio.png}{(e) Swin-UMamba}{0.170\textheight}
\hfill
\OriginalMambaPanel{figures/nnMamba_corrected.pdf}{(f) nnMamba}{0.170\textheight}
\par\vspace{2.0mm}
\OriginalMambaPanel{figures/HybridMamdrawio.png}{(g) HybridMamba}{0.170\textheight}
\hfill
\OriginalMambaPanel{figures/LKM UNet.png}{(h) LKM-UNet}{0.170\textheight}
\par\vspace{1.0mm}
\caption{Detailed state-space segmentation architectures complementing Figure~\ref{fig:mamba_architecture_progression}. (a) U-Mamba~\citep{ma2024umamba}; (b) SegMamba~\citep{xing2024segmamba}; (c) VM-UNet~\citep{ruan2024vm}; (d) Mamba-UNet~\citep{wang2024mamba}; (e) Swin-UMamba~\citep{liu2024swin}; (f) nnMamba~\citep{gong2025nnmamba}; (g) HybridMamba~\citep{wu2025hybridmamba}; and (h) LKM-UNet~\citep{wang2024lkm}. The panels expose differences in local feature extraction, spatial ordering, hierarchical processing, and fusion. Panel (f) summarizes nnMamba's Mamba-in-Convolution with Channel--Spatial Siamese learning (MICCSS). Experimental comparisons remain tied to the source-specific protocols in Table~\ref{tab:mamba_seg}. Panel (f) was revised using vector drawing code prepared with assistance from ChatGPT (OpenAI; historical model version not recorded).}
\label{fig:original_mamba_panels}
\end{figure*}
% END AUTHOR FIGURE fig:original_mamba_panels

\subsection{Selective State Propagation as an Alternative to Pairwise Attention}

A selective SSM propagates information through a latent state,
\begin{equation}
\mathbf{h}_t=\bar{\mathbf{A}}_t\mathbf{h}_{t-1}+\bar{\mathbf{B}}_t\mathbf{u}_t,
\qquad
\mathbf{y}_t=\mathbf{C}_t\mathbf{h}_t,
\label{eq:mamba_recurrence}
\end{equation}
where $\mathbf{u}_t\in\mathbb{R}^{d_{\mathrm{in}}}$ is the input token, $\mathbf{h}_t\in\mathbb{R}^{d_s}$ is the latent state, and $\mathbf{y}_t\in\mathbb{R}^{d_{\mathrm{out}}}$ is the output. The barred matrices $\bar{\mathbf{A}}_t$ and $\bar{\mathbf{B}}_t$ denote the discretized state-transition and input-projection operators after the input-dependent selective parameterization; $\mathbf{C}_t$ is the output projection. This is a compact recurrence used to discuss spatial serialization, not a second independent definition of the continuous-time SSM. Mamba's selective mechanism therefore provides input-dependent long-range propagation with linear complexity in sequence length~\citep{gu2023mamba}. In segmentation, however, the sequence is not given naturally: a mapping from a spatial grid or volume to an ordered token stream must first be chosen. Figure~\ref{fig:uneteval}(e) illustrates the generic SSM/Mamba stage in the broader architectural progression; the methods reviewed below differ mainly in how they construct this sequence, combine it with spatial operators, and restore dense predictions.

\subsection{Architectural Adaptation Strategies}

The literature can be read as four complementary design responses to the spatial-serialization problem. The first inserts SSM blocks into established convolutional pipelines; the second replaces larger portions of the network with visual SSM blocks; the third redesigns scanning for volumetric data; and the fourth adds channel--spatial processing, multi-scale interaction, or frequency-domain pathways when a single spatial scan is insufficient. This organization avoids treating every model as an isolated architecture and exposes the actual design variable shared by the family.

\subsubsection{From SSM Insertion to SSM-Dominant U-Shaped Models}

U-Mamba~\citep{ma2024umamba} takes the conservative route: convolution retains local feature extraction while SSM blocks augment the encoder or bottleneck with long-range propagation. nnMamba~\citep{gong2025nnmamba} moves beyond generic flattening by combining convolution with channel--spatial Siamese state-space processing through its MICCSS block. Both methods retain the stabilizing inductive structure of CNN-based medical segmentation while changing how context is propagated.

VM-UNet~\citep{ruan2024vm} shifts further toward an SSM-dominant design, using visual state-space blocks and multi-directional cross-scanning within a U-shaped architecture. Mamba-UNet~\citep{wang2024mamba} retains a more hybrid reconstruction path, combining a VMamba-style encoder with conventional skip transfer and convolutional decoding. UltraLight VM-UNet~\citep{wu2025ultralight} demonstrates the opposite end of the design spectrum by aggressively reducing parameters and computation for 2D segmentation, whereas SA-UMamba~\citep{liu2025sa} reintroduces position-sensitive local processing through a spatially enhanced state-space block. Read together, these methods show that replacing attention with an SSM does not eliminate the need for local spatial bias; it changes where and how that bias must be reintroduced.

\subsubsection{Pretraining and Hierarchical Hybrids}

Swin-UMamba~\citep{liu2024swin} couples hierarchical visual state-space modeling with pretrained initialization. Its importance is conceptual as much as architectural: performance gains cannot always be attributed to the state equation or scanning pattern alone because initialization quality and representation transfer may dominate small architectural changes. This observation connects the SSM literature to the foundation-model literature that follows: once representation quality becomes a major performance driver, architecture design and pretraining strategy can no longer be analyzed independently.

\subsubsection{From 2D Cross-Scanning to Volumetric Organization}

The decisive challenge for medical SSMs appears in 3D. SegMamba~\citep{xing2024segmamba} and SegMamba-V2~\citep{xing2025segmamba} use tri-oriented processing to expose axial, sagittal, and coronal relationships rather than relying on a single flattened order. nnMamba~\citep{gong2025nnmamba} uses channel--spatial Siamese processing within convolutional features, while other volumetric methods use plane-wise or hierarchical decompositions. These approaches provide different paths for inter-slice interaction. A shared constraint concerns ordering: a single linear ordering cannot make every 3D grid edge an adjacent sequence pair; whether scans, recurrent states, and local operators preserve task-relevant spatial information must be tested.

\subsubsection{Multi-Scale and Dual-Domain Extensions}

Specialized architectures broaden the role of scanning itself. LKM-UNet~\citep{wang2024lkm} combines pixel-level and patch-level state-space processing to cover fine and coarse dependencies. T-Mamba~\citep{hao2026t} introduces a frequency-aware pathway for boundary-sensitive dental segmentation, while HybridMamba~\citep{wu2025hybridmamba} integrates slice-oriented, local-adaptive, and frequency-gated processing. These methods are important because they do not treat ``global context'' as a single homogeneous objective: anatomical layout, local boundary structure, and spectral detail are routed through different representational paths.


\subsection{Architecture-Centric Comparison}


Mamba-based methods apply sequential state propagation to medical spatial structure as detailed in Table~\ref{tab:mamba_seg}. The corresponding CNN and transformer baselines can be found in Tables~\ref{tab:unet_compare} and \ref{tab:transformer_medseg} and are not reproduced here.

\begin{table*}[!t]
\centering
\fontsize{7.5}{9}\selectfont
\renewcommand{\arraystretch}{1.05}
\setlength{\tabcolsep}{3pt}
\begin{tabularx}{\textwidth}{>{\raggedright\arraybackslash}p{2.4cm}>{\raggedright\arraybackslash}p{1.5cm}>{\raggedright\arraybackslash}p{1.8cm}>{\raggedright\arraybackslash}p{1.1cm}>{\raggedright\arraybackslash}p{1.4cm}>{\raggedright\arraybackslash}p{1.0cm}>{\raggedright\arraybackslash}p{.95cm}>{\raggedright\arraybackslash}X}
\toprule
\textbf{Method / Variant} & \textbf{Mechanism} & \textbf{Benchmark} & \textbf{DSC (\%)} & \textbf{Boundary Metric} & \textbf{Params (M)} & \textbf{FLOPs (G)} & \textbf{Numerical Source / Config} \\
\midrule

U-MambaBot \citep{ma2024umamba}$^{\dagger}$
& CNN + SSM
& Synapse (8 organs)
& 86.88
& HD95 7.35
& 41.95
& 156.32
& Insertion of SSM in bottleneck. Numerical source: \citet{lumetti2025taming}, Table 3 and 5. Synapse-specific model and costs. \\

U-MambaEnc \citep{ma2024umamba}$^{\dagger}$
& CNN + SSM
& Synapse (8 organs)
& 87.82
& HD95 7.83
& 42.85
& 231.18
& Insertion of SSM in encoder. Same later source and benchmark as previous row. No cross-paper ranking with these means. \\

nnMamba \citep{gong2025nnmamba,review95_gong2024nnmamba}
& CNN + SSM, 3D
& BraTS 2023
& 89.97
& HD95 6.53
& NR
& NR
& Preprint v2, Table 1. MICCSS is a combination of channel/spatial Siamese processing and convolution. Table 2 15.55M/141.14G costs are for AMOS, not transferred here. \\

VM-UNet \citep{ruan2024vm}
& VSS encoder-decoder
& Synapse (8 organs)
& 81.08
& HD95 19.21
& NR
& NR
& Primary Table II. Four-direction spatial scanning with an asymmetric U-shaped model. Includes embedding/reconstruction cost. \\

Mamba-UNet \citep{wang2024mamba}
& VMamba U-shape
& Synapse CT
& 64.29
& HD 24.4725
& NR
& NR
& Primary Table 2 shows Dice 0.6429, HD 24.4725. This is own training protocol. Not comparable to VM-UNet protocol. \\

UltraLight VM-UNet \citep{wu2025ultralight}
& Parallel VSS
& ISIC2017 / ISIC2018 / PH2
& 90.91 / 89.40 / 92.65
& NR
& 0.049
& 0.060
& Published primary Table 1. Three Dice values are corresponding to the datasets. Model size and GFLOPs as reported there. No unmatched backbone ratio inferred. \\

SA-UMamba \citep{liu2025sa}
& VSS + spatial attention
& Synapse CT
& 82.54
& HD95 16.80
& 43.45
& 7.02
& Residual spatial state-space block. Combination of VSS and receptive-field attention convolution. Scores and costs from primary Tables 7-8. RSSS block schedule in encoder is $\{2,2,2,2\}$. In decoder $\{2,2,2,1\}$. \\

Swin-UMamba (Mamba decoder) \citep{liu2024swin,review95_liu2024swinumamba}
& 2D VSS hierarchy
& AbdomenMRI / Endoscopy
& 77.05 / 67.83
& NSD 0.8376 / 0.6933
& 28
& 18.9 / 45.3
& Preprint v2, Tables 1-2. Mamba decoder variant; NSD is a surface-overlap score on a 0--1 scale. ImageNet-initialized VMamba encoder. Different setup for 60M CNN decoder base. Both costs are 2D and dataset-specific. \\

SegMamba \citep{xing2024segmamba}$^{\dagger}$
& 3D tri-oriented SSM
& BraTS 2023
& 91.32
& HD95 3.56
& NR
& NR
& Tri-oriented spatial processing. Convolutional reconstruction. Numerical record: SegMamba baseline in HybridMamba \citep{wu2025hybridmamba}, Table 1. $128^3$ crops. Split in the same paper. \\

SegMamba-V2 \citep{xing2025segmamba}
& 3D tri-oriented SSM
& Volumetric benchmarks
& NV
& NV
& NV
& NV
& Hierarchical scale downsampling. Tri-oriented spatial Mamba. Design verified in primary abstract. Scores and their dispersion are not given until checking full experimental record. \\

T-Mamba \citep{hao2026t,review95_hao2024tmamba}
& Frequency + SSM
& CBCT tooth segmentation
& 93.60
& HD 1.18
& 1.04
& 25.43
& Expanded preprint v2; main 3D comparison. A learned gate selects frequency and spatial features. Table 2 reports PMFSNet3D at 91.30 DSC / HD 5.57 mm and T-Mamba at 93.60 / 1.18 mm; +2.30 Dice points and $-4.39$ mm HD are comparator-relative changes, not standard deviations. Most baseline results are inherited from the cited PMFSNet3D study. \\

LKM-UNet \citep{wang2024lkm}
& Pixel + patch SSM
& 3D Abdomen CT / 2D Abdomen MR
& 86.82 / 77.35
& NSD 90.02\% / 83.80\%
& NR
& NR
& Primary MICCAI Table 1; NSD is expressed as a percentage. Pixel and patch Mamba pathways process different spatial scales. Two scores are for different dimensionality and datasets. \\

HybridMamba \citep{wu2025hybridmamba}
& Spatial + frequency SSM
& BraTS 2023 / LC
& 91.92 / 75.34
& HD95 3.48 / 44.52
& NR
& NR
& Primary MICCAI Table 1. Pairs of values are corresponding to the datasets. Slice-oriented and local-adaptive Mamba. Frequency gating. $128^3$ crops. 70/10/20 split. \\

\bottomrule
\end{tabularx}
\caption{Representative Mamba-based segmentation architectures. Model variants, dataset-specific measurements, and numerical sources are listed per row. HD is not quietly converted to HD95. No assumption of common distance units. No pooling/ordering results in a leaderboard. NR, NV, and the provenance marker are defined in Table~\ref{tab:transformer_medseg}.}
\label{tab:mamba_seg}
\end{table*}

\paragraph{A naming caveat}
The MedMamba entry refers to \citet{wang2026medmamba}, \emph{MedMamba: Multi-scale deformable attention via state space models for robust medical image segmentation}, published in \emph{Biomedical Signal Processing and Control} 112, article 108363 (2026), DOI \href{https://doi.org/10.1016/j.bspc.2025.108363}{10.1016/j.bspc.2025.108363}. The publisher record identifies a medical image segmentation method, distinct from \emph{MedMamba: Vision Mamba for Medical Image Classification}. We retain the segmentation paper as architectural context only; no numerical comparison row is included because its experimental protocols and outcomes were not subjected to the configuration-level audit used for the tables.

\subsection{Provisional Mechanistic Synthesis: The Scanning Bottleneck}


The SSM/Mamba sample is purposive. We do not estimate the prevalence or maturity of mechanisms in the field. Statements below are \emph{transition-specific hypotheses}. They are not family-wide laws or claims of superiority. The inventory explicitly states coverage. Diagnostic tests indicate what additional evidence would support or contradict each explanation.

First, \textbf{scanning strategy is a key structural variable}. Four-direction cross-scanning, tri-oriented volumetric decomposition, channel--spatial state processing, pixel--patch dual-scale scanning, and slice-oriented processing all modify the neighborhood graph presented to the same underlying sequential mechanism. We should compare the resulting interaction paths using neighborhood distortion, coordinate equivariance, task performance, and realized cost, rather than asserting spatial neutrality.

Second, \textbf{linear asymptotic complexity does not imply universal deployment efficiency}. Lightweight 2D variants show large savings. But volumetric pipelines can have large constants from multiple directional scans, multiple planes, multi-scale branches, or heavy convolutional backbones. Thus, computational claims must be interpreted with dimensionality, input size, scan multiplicity, hardware implementation, etc., rather than just the $\mathcal{O}(N)$ label.

Third, \textbf{the 2D--3D representation problem remains open within the sampled literature}. Volumetric interaction paths are represented explicitly through plane-wise and directional decompositions, but preservation of task-relevant topology, orientation behavior, and neighborhood information requires targeted measurement. This motivates, rather than empirically validates, the proposed research direction on native 3D SSM design: spatial propagation that respects volumetric structure while retaining the computational advantages of state-space processing.

The distinction from foundation models is conceptual rather than chronological. Mamba changes \emph{how} context is propagated, whereas SAM changes \emph{where segmentation knowledge comes from and how predictions are conditioned}.


\section{SAM-Based Medical Segmentation: From Promptable Transfer to Clinical Adaptation}
\label{sec:sam_review}


The shift from CNNs to Transformers and SSMs primarily changes the mechanism of representation formation and contextual interaction. SAM introduces another axis of change. Segmentation is conditioned on a pretrained representation and external prompt, rather than being learned solely as a fixed task-specific mapping. A promptable segmenter can be written as
\begin{equation}
\hat{Y}=f_{\theta}(X,P),
\end{equation}
where $X$ is the medical image, $P$ is a spatial, semantic, textual, or propagated prompt, and $\hat{Y}$ is the predicted mask. Thus, the medical adaptation problem has two coupled sources of uncertainty. Is the pretrained image representation aligned with the clinical domain? Is the conditioning signal informative enough?

Rather than listing models, we organize the literature around five adaptation questions. How can we align a natural-image foundation model with medical data? How can we specialize it without full retraining? How can we introduce volumetric continuity? How can we reduce the reliance on manual prompts? How can we extend prompt conditioning to temporal memory or clinical semantics? These questions map directly to the architectural variants shown in Fig.~\ref{fig:sam_medical_architecture_progression} and compared in Table~\ref{tab:detailed_comparison}.


\begin{figure*}[!t]
\centering
\includegraphics[width=0.8\textwidth]{assets/sam_adaptation_mechanisms.pdf}
\caption{Four adaptation mechanisms for promptable medical segmentation. (a) SAM baseline. Mask decoder combines image and prompt embeddings. Zero-shot use shown (no fine-tuning for medical domain)~\citep{kirillov2023sam}. (b) Medical supervision can update selected pretrained weights or parameter-efficient modules. E.g., MedSAM, SAMed, Med-SA. Branches indicate alternative strategies rather than a single joint architecture~\citep{ma2024medsam,zhang2023samed,wu2025medsa}. (c) Volumetric image/mask processing in SAM-Med3D, SegVol. Supported prompt types and inference stages are method-specific~\citep{wang2025sam,du2024segvol}. (d) Two different mechanisms to reduce repeated interaction. MedLSAM localizes anatomy to generate prompts~\citep{lei2025medlsam}. Medical SAM 2 uses stored context between sequential inference steps~\citep{zhu2024medical}. Memory feedback indicates use in the next step. These are original schematics showing adaptation choices. Not identical layers or a ranking of segmentation accuracy. Drawing code and layout were prepared with assistance from ChatGPT (OpenAI; historical model version not recorded).}
\label{fig:sam_medical_architecture_progression}
\end{figure*}

% BEGIN AUTHOR FIGURE fig:original_sam_panels
Figure~\ref{fig:original_sam_panels} supplements the summary of mechanisms with architectural views of the medical SAM variants we discuss below.

\begin{figure*}[!t]
\centering
\setlength{\tabcolsep}{0pt}
\newcommand{\OriginalSAMPanel}[3]{%
  \begin{minipage}[t]{0.485\textwidth}
    \centering
    \begin{minipage}[c][#3][c]{\linewidth}
      \centering
      \includegraphics[width=\linewidth,height=#3,keepaspectratio]{#1}
    \end{minipage}\par\vspace{1.2mm}
    {\footnotesize\bfseries #2\par}
  \end{minipage}%
}
\OriginalSAMPanel{figures/SAM_Med3ddrawio.png}{(a) SAM-Med3D}{0.220\textheight}
\hfill
\OriginalSAMPanel{figures/MedSAMFdrawio.png}{(b) MedSAM}{0.220\textheight}
\par\vspace{2.0mm}
\OriginalSAMPanel{figures/SAMed_LoRAF.drawio.png}{(c) SAMed}{0.220\textheight}
\hfill
\OriginalSAMPanel{figures/MedSAM2_corrected.pdf}{(d) Medical SAM 2}{0.220\textheight}
\par\vspace{2.0mm}
\OriginalSAMPanel{figures/MEdSAdrawio.png}{(e) Med-SA}{0.220\textheight}
\hfill
\OriginalSAMPanel{figures/SegVolFdrawio.png}{(f) SegVol}{0.220\textheight}
\par\vspace{1.0mm}
\caption{Detailed promptable segmentation architectures complementing Figure~\ref{fig:sam_medical_architecture_progression}. (a) SAM-Med3D~\citep{wang2025sam}. (b) MedSAM~\citep{ma2024medsam}. (c) SAMed~\citep{zhang2023samed}. (d) Medical SAM 2~\citep{zhu2024medical}. (e) Med-SA~\citep{wu2025medsa}. (f) SegVol~\citep{du2024segvol}. The panels illustrate volumetric prompting, medical fine-tuning, low-rank adaptation, sequential memory, adapter-based specialization, and interactive volumetric segmentation. Trainable modules, supported prompt types, and inference protocols vary across methods. Numerical results are separately listed in Table~\ref{tab:detailed_comparison}. Panel (d) was revised using vector drawing code prepared with assistance from ChatGPT (OpenAI; historical model version not recorded).}
\label{fig:original_sam_panels}
\end{figure*}
% END AUTHOR FIGURE fig:original_sam_panels
\subsection{From Zero-Shot Mismatch to Domain Alignment}

Vanilla SAM~\citep{kirillov2023sam} provides the promptable baseline, but its natural-image pretraining does not necessarily transfer to medical scans, where boundaries may be low contrast, targets small or diffuse, and intensity distributions modality-specific. Anatomical identity may also depend on spatial context rather than appearance alone. Adaptation can modify the image representation, spatial interaction, prompt interface, or mask decoder; the appropriate intervention depends on the target task.

MedSAM~\citep{ma2024medsam} addresses this mismatch through medical-domain fine-tuning across heterogeneous modalities, while MA-SAM~\citep{chen2024ma} investigates modality-agnostic adaptation. Other strategies use parameter-efficient updates or spatial adapters to extend two-dimensional representations to volumetric tasks. Task-specific systems such as PolypSAM~\citep{li2024polyp} and SkinSAM~\citep{hu2024skinsam} illustrate further specialization. The central question is how much adaptation is needed for clinically useful performance and whether it preserves transfer beyond the target dataset. Neither extensive fine-tuning nor a strong in-domain score alone establishes retained generality.

Related vision--language methods extend medical adaptation beyond the SAM family. MedCLIPSeg~\citep{koleilat2026medclipseg_cvpr} uses probabilistic cross-modal attention over image and text features and a soft contrastive objective to study data-efficient and generalizable segmentation. Its evaluation addresses a different conditioning regime and is not interpreted as a prompt-matched comparison with the SAM records in Table~\ref{tab:detailed_comparison}.

\subsection{Parameter-Efficient Specialization and Representation Preservation}

As an alternative to full fine-tuning, much of the pretrained backbone can be frozen while a small task-specific update is learned. SAMed~\citep{zhang2023samed} and LoRA-MedSAM~\citep{hu2024lora} use low-rank updates, whereas Med-SA~\citep{wu2025medsa} and SAM-Med2D~\citep{cheng2023sam} introduce adapter-style modules. AdaptiveSAM~\citep{paranjape2024adaptivesam} explores lighter adaptation options. These methods differ in the parameters updated, conditioning mechanism, and resource demands, which must be reported when comparing them.

The importance of parameter-efficient adaptation extends beyond the number of trainable parameters: it separates a shared pretrained representation from the task-specific correction. This can be advantageous for multi-domain systems, but a restricted update may not accommodate a large shift in modality or target anatomy. I-MedSAM~\citep{wei2024imedsam} explores an implicit neural-representation decoder for continuous boundary prediction, while distillation and self-prompting address different efficiency or automation needs. MoE-SAM~\citep{li2025moe} introduces sparse expert specialization for heterogeneous feature distributions. Whether these interventions retain broader transfer requires evaluation outside their adaptation settings.

\subsection{From 2D Prompting to Volumetric Continuity}

Domain alignment alone does not address volumetric consistency in CT or MRI. Adaptation generally follows two routes: \textbf{architectural reconstruction}, using native or near-native 3D promptable systems such as SAM-Med3D~\citep{wang2025sam} and SegVol~\citep{du2024segvol}; or \textbf{bridging adaptation}, which introduces depth-aware processing into a predominantly 2D pretrained model, as in 3DSAM-adapter~\citep{gong20243dsam} and the space--depth components of Med-SA~\citep{wu2025medsa}. The relevant comparison is not merely 2D versus 3D, but the source of inter-slice information and its effect on prediction consistency.

Volumetric adaptation introduces inter-slice processing, but its effects on accuracy, data requirements, retained pretraining benefits, and deployment cost depend on the implementation. Lightweight bridges can preserve a pretrained 2D backbone while adding spatial interactions; fully volumetric designs alter a larger portion of the pipeline. SAM-Med3D-MoE~\citep{wang2024sam} adds routing between general and task-specific decoders, illustrating that volumetric continuity and specialization are separate design dimensions. Matched experiments are needed to distinguish their contributions.

Free-text conditioning introduces another form of volumetric interaction. VoxTell~\citep{rokuss2026voxtell_cvpr} combines multi-stage visual--language fusion with volumetric mask prediction from natural-language descriptions. This makes prompt specificity, linguistic variation, and transfer to unseen concepts relevant evaluation conditions, distinct from the point- and box-prompt protocols tabulated below.

\subsection{From Manual Interaction to Operational Autonomy}

Promptability is a defining feature of SAM, but repeated manual prompts can limit deployment at scale. MedLSAM~\citep{lei2025medlsam} addresses this by localizing anatomy to generate prompts for 3D CT. Such automation may reduce user interaction, although the resulting segmentation remains dependent on localization quality. Prompt-generation and segmentation errors should therefore be evaluated separately.

Other methods modify the conditioning mechanism itself. DeSAM~\citep{gao2024desam} separates prompt-relevance estimation from dense mask reconstruction, while SAM-MPA~\citep{xu2024sammpa} and SSM-SAM~\citep{leng2024self} investigate conditioning from limited support information. ScribblePrompt-SAM~\citep{wong2024scribbleprompt} broadens the interaction interface, SurgicalSAM~\citep{yue2024surgicalsam} uses class-conditioned prototypes, and MedSeg-R~\citep{shao2025medseg} incorporates clinically structured semantic priors. Together these approaches shift conditioning from manually specified coordinates toward automatically generated, propagated, or semantically derived information, with different failure modes for each source.

\subsection{Temporal Memory and Sequential Prompt Propagation}

SAM2-based adaptations introduce sequential memory into medical segmentation. Medical SAM 2~\citep{zhu2024medical} reuses information from preceding frames or slices to condition subsequent predictions. This may reduce repeated prompting and improve continuity, but also permits errors in earlier masks or stored context to propagate through the sequence. The potential reduction in annotation effort must therefore be assessed alongside memory reliability and error accumulation.

Figure~\ref{fig:sam_medical_architecture_progression} summarizes representative adaptation mechanisms, including backbone updating, volumetric processing, prompt generation, and sequential memory. These panels provide architectural context; source-specific performance values and inference protocols are reported separately in Table~\ref{tab:detailed_comparison}.


Related methods extend this design space in different ways. SAM-UNETR investigates transfer for clinically significant prostate cancer; SAM2-UNet uses a SAM 2 encoder in a U-shaped segmentation design; and SAM-TransUNet incorporates automatic cues in transformer-based segmentation~\citep{alzate2023sam,xiong2026sam2,qi2025sam}. MedicoSAM and Rep-MedSAM further investigate medical adaptation and deployment efficiency~\citep{archit2025medicosam,wei2024rep}. These studies are included as architectural context, not as a basis for ranking results obtained under different protocols.

\subsection{Consolidated Comparison}


Table~\ref{tab:detailed_comparison} consolidates source-specific outcomes alongside dataset, metric, prompting, and inference conditions. The adaptation and architectural-difference columns identify the trainable scope and the relevant design mechanism; fields that could not be verified remain explicitly unresolved. The following synthesis therefore concerns mechanism-level implications rather than a pooled numerical ranking.


\begin{table*}[p]
\centering
\fontsize{7.5}{9}\selectfont
\renewcommand{\arraystretch}{1.12}
\setlength{\tabcolsep}{3.0pt}

\begin{tabularx}{\textwidth}{>{\raggedright\arraybackslash}p{2.0cm}>{\raggedright\arraybackslash}p{2.0cm}>{\raggedright\arraybackslash}p{2.35cm}>{\raggedright\arraybackslash}p{1.2cm}>{\raggedright\arraybackslash}p{1.7cm}>{\raggedright\arraybackslash}p{1.4cm}>{\raggedright\arraybackslash}X}
\toprule
\textbf{Method}
& \textbf{Dataset / Target}
& \textbf{Prompt or Evaluation Protocol}
& \textbf{Metric}
& \textbf{Reported Value}
& \textbf{Adaptation}
& \textbf{Architectural / Methodological Difference} \\
\midrule

SAM \citep{kirillov2023sam}
& BTCV, 12 organs
& Three point prompts per slice. Re-evaluated in Med-SA study
& Mean DSC
& 63.1\%\textsuperscript{a}
& Zero-shot
& Original SAM. Frozen natural image representation. No medical-domain adaptation. \\

MedSAM \citep{ma2024medsam}
& BTCV, 12 organs
& Three point prompts per slice. Re-evaluated in Med-SA study
& Mean DSC
& 82.0\%\textsuperscript{a}
& Study-specific full tuning
& Med-SA evaluation lists all 636M parameters as tunable for this baseline. This is the evaluated implementation. Not a claim about all MedSAM training protocols. \\

Med-SA \citep{wu2025medsa}
& BTCV, 12 organs
& Three point prompts per slice
& Mean DSC
& 88.7\%\textsuperscript{a}
& Adapter
& Freezes most SAM parameters. Introduces SD-Trans for depth-aware processing. HyP-Adpt for prompt-conditioned adaptation. \\

Med-SA \citep{wu2025medsa}
& BTCV, 12 organs
& Per-slice bounding box. 0.75 box-quality setting
& Mean DSC
& 89.8\%\textsuperscript{a}
& Adapter
& Best reported BTCV configuration in Med-SA study. 13M of 636M parameters updated. \\

SAMed \citep{zhang2023samed}
& Synapse, 8 organs
& Automatic multi-class semantic segmentation
& DSC / HD
& 81.88\% / 20.64
& LoRA
& Inserts low-rank updates in ViT image encoder. Fine-tunes prompt encoder and mask decoder. \\

SAM-Med2D \citep{cheng2023sam}
& AMOS22, 15 organs
& Point prompts. Five-step corrective interaction. Re-evaluated in SegVol study
& Mean of category median DSC
& 66.68\%\textsuperscript{b}
& 2D adapter
& Large-scale two-dimensional medical data. Adapter-based tuning. Keeps slice-wise processing. \\

SAM-Med3D \citep{wang2025sam}
& AMOS-CT
& One 3D point per target
& Mean DSC
& 79.94\%
& Native 3D
& Fully learnable 3D image, prompt, mask-decoder components. Trained on SA-Med3D-140K corpus. \\

SAM-Med3D \citep{wang2025sam}
& AMOS-CT
& Ten 3D corrective points per target
& Mean DSC
& 83.99\%
& Native 3D
& Shows the improvement from additional volumetric interaction. Same benchmark. \\

SegVol \citep{du2024segvol}
& AMOS22, 15 organs
& Bounding-box and text prompts
& Mean of category median DSC
& 85.93\%\textsuperscript{b}
& Native 3D
& Spatial and semantic prompts. Zoom-out--zoom-in inference mechanism for volumetric segmentation. \\

3DSAM-adapter \citep{gong20243dsam}
& KiTS21 kidney tumor
& One point per volume
& DSC / NSD
& 73.78\% / 83.86\%
& 3D adapter
& Holistically adapts SAM from 2D to 3D. Keeps most pretrained parameters. Lightweight spatial adapters. \\

3DSAM-adapter \citep{gong20243dsam}
& MSD pancreas tumor
& One point per volume
& DSC / NSD
& 54.09\% / 76.27\%
& 3D adapter
& Same volumetric prompt protocol. Small, low-contrast tumor segmentation task. \\

SAM-Med3D-MoE \citep{wang2024sam}
& SPPIN 2023 neuroblastoma
& Bounding-box prompt
& DSC
& 48.9\%
& MoE-3D
& Routes between general SAM-Med3D decoder and task-specific expert decoders. Corresponding general expert is 32.3\% DSC. \\

DeSAM \citep{gao2024desam}
& Six-site prostate segmentation
& DeSAM-P: $9\times9$ grid of point prompts
& Mean DSC across six sites
& 79.02\%
& Decoder decoupling
& Decouples prompt-relevant mask generation from prompt-invariant image representation. Improves cross-domain generalization. \\

Medical SAM 2 \citep{zhu2024medical,review95_zhu2024medical_v1}
& BTCV, 12 organs
& Preprint v1 3D protocol. Prompt probability 0.25
& Mean DSC
& 88.57\%
& SAM2 memory
& Treats medical slices as video frames. Propagates predictions using confidence-aware memory bank. Weighted memory retrieval. \\

SurgicalSAM \citep{yue2024surgicalsam}
& EndoVis2018
& Prototype-derived class prompts. No point/box detector
& Challenge IoU /
mcIoU
& 80.33\% /
58.87\%
& Class-prompt tuning
& Prototype-based class-prompt encoder. Contrastive prototype learning. End-to-end surgical instrument segmentation. \\

\bottomrule
\end{tabularx}

\caption{Protocol-annotated reported comparison of representative SAM-based medical segmentation methods. Every value is associated with a specific dataset, metric, prompt/inference protocol. Values with \textsuperscript{a} are reproduced from Table 2 of the published Med-SA study \citep{wu2025medsa}. Prompt budget is per slice. These SAM/MedSAM entries are the evaluated baselines in that study. Values with \textsuperscript{b} are from Table 2 of the NeurIPS 2024 SegVol study \citep{du2024segvol}. Per-category values are median Dice. Displayed overall value is average of 15 categories. Multiplying by 100 only changes displayed scale. SAM-Med2D has five corrective interactions. SegVol has box and text prompts. So, these two rows are not a matched-prompt comparison. SAM-Med3D values from Table 4 of arXiv:2310.15161v3. 3DSAM-adapter values from Table 1 of author preprint arXiv:2306.13465. DeSAM values from its MICCAI Table 2. 79.02 is only prostate aggregate. Medical SAM 2 values from arXiv:2408.00874v1, Table 1. SurgicalSAM uses the main paper EndoVis2018 setting. Not the average of additional seeds. Results obtained with different datasets, prompts, interaction budgets, dimensionality, evaluation procedures are not directly comparable. Do not interpret as a unified leaderboard. DSC = Dice similarity coefficient. HD = Hausdorff distance as reported by the source. NSD = normalized surface Dice. mcIoU = mean class intersection over union. MoE = mixture of experts.}
\label{tab:detailed_comparison}
\end{table*}
\paragraph{Constraint redistribution: mechanistic perspective on adaptation}

SAM-based medical segmentation involves four interconnected considerations: \textit{representation quality}, \textit{medical-domain alignment}, \textit{volumetric continuity}, and \textit{operational autonomy}. Fine-tuning can improve target-task fit, but its resource requirements and effects on transfer depend on the update strategy and data. Parameter-efficient tuning limits changes to the pretrained backbone but may restrict the range of domain shifts it can accommodate. Volumetric processing introduces inter-slice information whose effect on consistency and cost should be measured. Automatic or semantic prompting reduces manual interaction but creates dependencies on localizers, support examples, semantic priors, or sequential memory.

The same bounded question applies to these adaptations: which limitation is reduced, and which dependency, if any, becomes important under the declared task and evaluation protocol? Candidate dependencies include specialization, conditioning reliability, dimensionality, and computational cost.

The clinical implications depend on the intended workflow. Promptable models may be useful for interactive annotation, few-shot adaptation, repeated tasks, or sequential propagation. Fully automated deployment additionally requires assessment of target-task accuracy, conditioning reliability, external-domain performance, and failure handling; volumetric applications require tests of spatial consistency. These evaluation requirements link the SAM literature to the cross-paradigm analysis below.



\subsection{Targeted Contextual Update Through November 2026}
\label{subsec:context_update_2026}

A targeted update checked on 20 November 2026 adds architectural and evaluation context from seven CVPR 2026 and sixteen MICCAI 2026 publications, including the three CVPR studies already discussed above. These 23 publications are a separate contextual extension to the retrospective synthesis. They do not change its 30 June 2026 boundary, the frozen 95-family inventory, or the configuration-specific numerical evidence in S2. The following descriptions summarize the proposed mechanisms; they do not establish matched improvements across papers.

\paragraph{Local detail and state propagation}
HEDS-Net combines state-space processing with an axial skip bridge and progressively weighted supervision, whereas VM-NeXT UNet uses parallel convolutional and visual-state-space encoders with gated skips~\citep{fan2026hedsnet_context,cao2026vmnext_context}. AFMamba introduces frequency-guided parallel state-space processing for polyps, and BDFR-Net combines frequency and spatial features with skip alignment for fetal ultrasound~\citep{jiang2026afmamba_context,zhang2026bdfrnet_context}. These designs motivate separate ablations of propagation, boundary cues, and feature transfer under a common training budget.

\paragraph{Volumetric interaction}
DEM-Mamba uses directional wavelet features and balanced bidirectional scanning, while SRMA-Mamba combines anatomical multi-plane scans with reverse-attention refinement~\citep{zhu2026demmamba_context,zeng2026srmamamba_context}. TranSamba adds cross-plane Mamba blocks to a transformer under plane-level supervision~\citep{lyu2026transamba_context}. PlaneCycle instead reuses a two-dimensional pretrained backbone through cyclic aggregation over orthogonal planes~\citep{yu2026planecycle_context}. Its lifting operator requires no parameter updates, but downstream linear probing and fine-tuning remain distinct evaluation regimes. Voxel anisotropy, label availability, and full-volume inference cost are therefore necessary comparison conditions.

\paragraph{Pretrained representations and domain adaptation}
Dino U-Net couples a frozen DINOv3 encoder to dual-branch adaptation and fidelity-aware projection; Glasses combines DINOv3 adaptation with a boundary-sensitive convolutional branch and fine-tunes terminal backbone layers~\citep{li2026dinounet_context,zhu2026glasses_context}. SCFNet adapts convolutional aggregation and decoder fusion to the input, while SSUG combines spectral variation, uncertainty-guided gating, and consistency for SAM generalization~\citep{tang2026scfnet_context,li2026ssug_context}. APEX-SAM retrieves support exemplars and constructs anatomical prompts; spatio-semantic contextual prompting fuses spatial and semantic priors into dense prompts for frozen SAM-Med3D~\citep{mao2026apexsam_context,chattopadhyay2026contextprompts_context}. Their transfer claims require explicit support--query separation and untouched external domains.

\paragraph{Adaptation cost, supervision, and reasoning}
Dual-Adaptive SAM3 routes visual and textual inputs through low-rank experts, whereas IAD transfers teacher knowledge through semantic alignment and student-oriented assimilation~\citep{chen2026dasam3_context,hong2026iad_context}. WeakMed uses mask-to-box alignment and scale consistency during training; SD-FSMIS adapts diffusion features through support--query interaction and visual conditioning~\citep{wei2026weakmed_context,li2026sdfsmis_context}. CG-Reasoner links mask prediction with centroid-based positional reasoning~\citep{polamreddy2026cgreasoner_context}. These mechanisms should be assessed with explicit training, annotation, prompting, and inference budgets; language consistency is a separate outcome from mask accuracy.

\paragraph{Risk-controlled evaluation}
Conformal Lesion Segmentation calibrates a post-processing threshold to control case-level violations of a voxel-wise false-negative tolerance under exchangeability~\citep{tan2026cls_context}. It is evaluation and post-processing context rather than an additional architectural family. The guarantee does not independently control the false-positive rate or remove the need to examine distribution shift.

\section{Open challenges and future directions}
\label{sec:unified_challenges}


The architecture-specific analyses identify plausible links between information-flow mechanisms and task-dependent constraints. Convolutional context formation, attention connectivity, state propagation, pretrained conditioning, and feature fusion motivate distinct diagnostic tests. We organize these relationships as \emph{mechanism $\rightarrow$ intended benefit $\rightarrow$ prespecified constraint}. Published evidence and review-level interpretations remain separate, particularly when a controlled comparison has not yet been reported.

Figure~\ref{fig:cross_paradigm_summary} maps the identified constraints to the five research directions. The analysis proceeds from observed failure signatures and candidate mechanisms (Table~\ref{tab:failure_matrix}), to transition-specific evidence anchors (Table~\ref{tab:redistribution_evidence}), and then to protocol admissibility (Section~\ref{subsec:tradeoff_decomp}). The roadmap translates remaining uncertainties into prospective comparisons without treating heterogeneous published results as interchangeable measurements.

\begin{figure*}[!t]
\centering
\begingroup

\definecolor{ReviewFusion}{HTML}{2E7D32}
\definecolor{ReviewFusionFill}{HTML}{E5F3E3}
\definecolor{ReviewConcern}{HTML}{C00000}
\definecolor{ReviewConcernFill}{HTML}{FFF0F0}
\definecolor{ReviewLine}{HTML}{555555}

\begin{tikzpicture}[
    every node/.style={
        font=\sffamily\fontsize{7.6}{9}\selectfont,
        text=black
    },
    arr/.style={
        -{Latex[length=1.5mm]},
        line width=0.6pt,
        draw=ReviewLine
    },
    bott/.style={
        rounded corners=2pt,
        draw=ReviewConcern,
        fill=ReviewConcernFill,
        line width=0.6pt,
        text width=70mm,
        minimum height=12mm,
        align=flush left,
        inner xsep=1mm,
        inner ysep=0.6mm
    },
    chal/.style={
        rounded corners=2pt,
        draw=ReviewFusion,
        fill=ReviewFusionFill,
        line width=0.6pt,
        text width=70mm,
        minimum height=12mm,
        align=left,
        inner xsep=1mm,
        inner ysep=0.6mm
    },
    head/.style={
        font=\sffamily\fontsize{8.2}{9.5}\selectfont\bfseries,
        align=center,
        text width=72mm
    }
]

% Headings
\node[head] at (38mm,0)
    {Candidate architectural constraint};

\node[head] at (121mm,0)
    {Open challenge and objectives};

% Row 1
\node[bott] (b1) at (38mm,-11mm) {
    \textbf{Serialization and scan order.}
    SSM scans map 2D/3D neighborhoods to
    1D sequences (Sec.~\ref{sec:mamba_review}).
};

\node[chal] (c1) at (121mm,-11mm) {
    \textbf{OC1: Native 3D state-space modeling}\\
    Sec.~\ref{subsec:native3d}
    \quad \textit{GC, CE}
};

% Row 2
\node[bott] (b2) at (38mm,-24.5mm) {
    \textbf{Domain and initialization dependence.}
    Pretrained models may be sensitive to
    domain shifts (Sec.~\ref{sec:sam_review}).
};

\node[chal] (c2) at (121mm,-24.5mm) {
    \textbf{OC2: Medical-native foundation pretraining}\\
    Sec.~\ref{subsec:pretraining}
    \quad \textit{CG, GC}
};

% Row 3
\node[bott] (b3) at (38mm,-38mm) {
    \textbf{Feature selection and alignment.}
    Skip connections, hybrid fusion, and
    routing require compatible features.
};

\node[chal] (c3) at (121mm,-38mm) {
    \textbf{OC3: Spatial uncertainty-conditioned fusion}\\
    Sec.~\ref{subsec:adaptive_fusion}
    \quad \textit{LP, GC, CG}
};

% Row 4
\node[bott] (b4) at (38mm,-51.5mm) {
    \textbf{Boundary and high-frequency loss.}
    Downsampling and tokenization may
    suppress fine spatial detail.
};

\node[chal] (c4) at (121mm,-51.5mm) {
    \textbf{OC4: Frequency-aware boundary modeling}\\
    Sec.~\ref{subsec:frequency}
    \quad \textit{LP}
};

% Row 5
\node[bott] (b5) at (38mm,-65mm) {
    \textbf{Isolated evaluation objectives.}
    Single-metric gains may leave cost,
    boundary, and domain behavior untested.
};

\node[chal] (c5) at (121mm,-65mm) {
    \textbf{OC5: Multi-objective architecture selection}\\
    Sec.~\ref{subsec:moo}
    \quad \textit{LP, GC, CE, CG}
};

% Connections
\foreach \i in {1,...,5} {
    \draw[arr] (b\i.east) -- (c\i.west);
}

\end{tikzpicture}

\endgroup

\caption{
Mapping candidate architectural constraints to five
open challenges in medical image segmentation.
The analysis draws on
Sections~\ref{sec:cnn}--\ref{sec:sam_review}
and Tables~\ref{tab:failure_matrix}
and~\ref{tab:redistribution_evidence}.
Each challenge is linked to a testable hypothesis
and prespecified comparison
(Table~\ref{tab:research_roadmap}).
Objective abbreviations follow
Table~\ref{tab:construct_definitions}.
No new experiments are reported. Drawing code and layout were prepared with assistance from ChatGPT (OpenAI; historical model version not recorded).}

\label{fig:cross_paradigm_summary}
\end{figure*}


% \begin{figure*}[!t]
% \centering
% \begingroup

% \definecolor{ReviewFusion}{HTML}{2E7D32}
% \definecolor{ReviewFusionFill}{HTML}{E5F3E3}
% \definecolor{ReviewConcern}{HTML}{C00000}
% \definecolor{ReviewConcernFill}{HTML}{FFF0F0}
% \definecolor{ReviewLine}{HTML}{555555}

% \resizebox{\textwidth}{!}{%
% \begin{tikzpicture}[
%     every node/.style={
%         font=\sffamily\fontsize{8}{9.3}\selectfont,
%         text=black
%     },
%     arr/.style={
%         -{Latex[length=1.7mm]},
%         line width=0.6pt,
%         draw=ReviewLine
%     },
%     bott/.style={
%         rounded corners=2pt,
%         draw=ReviewConcern,
%         fill=ReviewConcernFill,
%         line width=0.6pt,
%         text width=71mm,
%         minimum height=16mm,
%         align=flush left,
%         inner xsep=1.7mm,
%         inner ysep=1.1mm
%     },
%     chal/.style={
%         rounded corners=2pt,
%         draw=ReviewFusion,
%         fill=ReviewFusionFill,
%         line width=0.6pt,
%         text width=71mm,
%         minimum height=16mm,
%         align=flush left,
%         inner xsep=1.7mm,
%         inner ysep=1.1mm
%     },
%     head/.style={
%         font=\sffamily\fontsize{8.5}{10}\selectfont\bfseries,
%         align=center,
%         text width=77mm
%     }
% ]

% % Column headings
% \node[head] at (42mm,0)
%     {Candidate architectural constraint};

% \node[head] at (131mm,0)
%     {Open challenge and objectives};

% % Row 1
% \node[bott] (b1) at (42mm,-12mm) {
%     \textbf{Serialization and scan order.}
%     SSM/Mamba scans approximate 2D/3D neighborhoods
%     through 1D sequences
%     (Sec.~\ref{sec:mamba_review}).
% };

% \node[chal] (c1) at (131mm,-12mm) {
%     \textbf{OC1: Native 3D state-space modeling}\\
%     Sec.~\ref{subsec:native3d}\\
%     \textit{Objectives: GC, CE}
% };

% % Row 2
% \node[bott] (b2) at (42mm,-29.5mm) {
%     \textbf{Domain and initialization dependence.}
%     Foundation-model adaptation and pretrained
%     hybrids introduce transfer sensitivity
%     (Secs.~\ref{sec:mamba_review}--\ref{sec:sam_review}).
% };

% \node[chal] (c2) at (131mm,-29.5mm) {
%     \textbf{OC2: Medical-native foundation pretraining}\\
%     Sec.~\ref{subsec:pretraining}\\
%     \textit{Objectives: CG, GC}
% };

% % Row 3
% \node[bott] (b3) at (42mm,-47mm) {
%     \textbf{Feature selection and alignment.}
%     Skip redesign, hybrid fusion, and routing
%     introduce feature correspondence challenges
%     (Secs.~\ref{sec:cnn}--\ref{sec:sam_review}).
% };

% \node[chal] (c3) at (131mm,-47mm) {
%     \textbf{OC3: Spatial uncertainty-conditioned fusion}\\
%     Sec.~\ref{subsec:adaptive_fusion}\\
%     \textit{Objectives: LP, GC, CG}
% };

% % Row 4
% \node[bott] (b4) at (42mm,-64.5mm) {
%     \textbf{Boundary and high-frequency loss.}
%     Downsampling, tokenization, and patch merging
%     may suppress thin structures and contours
%     (Table~\ref{tab:failure_matrix}).
% };

% \node[chal] (c4) at (131mm,-64.5mm) {
%     \textbf{OC4: Frequency-aware boundary modeling}\\
%     Sec.~\ref{subsec:frequency}\\
%     \textit{Objective: LP}
% };

% % Row 5
% \node[bott] (b5) at (42mm,-82mm) {
%     \textbf{Isolated evaluation objectives.}
%     Single-metric gains may leave efficiency,
%     boundary fidelity, and cross-domain
%     behavior untested
%     (Sec.~\ref{subsec:tradeoff_decomp}).
% };

% \node[chal] (c5) at (131mm,-82mm) {
%     \textbf{OC5: Multi-objective architecture selection}\\
%     Sec.~\ref{subsec:moo}\\
%     \textit{Objectives: LP, GC, CE, CG}
% };

% % Mapping arrows
% \foreach \i in {1,...,5} {
%     \draw[arr] (b\i.east) -- (c\i.west);
% }

% % Compact explanatory footer
% \node[
%     draw=ReviewLine,
%     rounded corners=2pt,
%     fill=white,
%     line width=0.5pt,
%     text width=159mm,
%     minimum height=9.5mm,
%     align=center,
%     inner sep=1.2mm,
%     anchor=north west,
%     font=\sffamily\fontsize{7.8}{9}\selectfont
% ] at (5mm,-94mm) {
%     Each challenge is associated with a testable
%     hypothesis and a prespecified diagnostic
%     comparison (Table~\ref{tab:research_roadmap}).
%     No new experiments are reported.
% };

% \end{tikzpicture}%
% }

% \endgroup

% \caption{
% Mapping candidate architectural constraints to five
% open challenges in medical image segmentation.
% The analysis draws on
% Sections~\ref{sec:cnn}--\ref{sec:sam_review}
% and Tables~\ref{tab:failure_matrix}
% and~\ref{tab:redistribution_evidence}.
% Objective abbreviations follow
% Table~\ref{tab:construct_definitions}.
% }

% \label{fig:cross_paradigm_summary}
% \end{figure*}

\subsection{Mechanism-linked failure modes: cross-paradigm analysis}
\label{subsec:failure_modes}

An observed benchmark difference does not by itself identify its cause. An architectural explanation gains support when errors arise on cases that stress the proposed mechanism, a targeted intervention changes those errors, and relevant data, initialization, optimization, and evaluation conditions are controlled. We therefore distinguish an \emph{observed error}, a \emph{candidate mechanism}, and \emph{intervention-based evidence}. Plausible information flow alone does not establish a structural failure.

The relevant mechanisms differ between the families reviewed here. Convolutional models build context through successive local operations; attention defines content-dependent interaction graphs; SSMs propagate compressed states through selected paths; and promptable models use pretrained representations with external conditioning. Skip connections and hybrid branches affect how these representations reach the decoder. Each mechanism suggests a diagnostic question, but its existence does not mean the corresponding limitation is active in a particular task.

\subsubsection{Failure matrix across paradigms}

Table~\ref{tab:failure_matrix} links candidate mechanisms to observable errors and discriminating interventions. It is important to separate the sensitivity of an already trained model from the effect of choosing an architectural design. Changing scan order at inference probes fixed-model sensitivity; evaluating the design itself requires separately trained variants under controlled conditions. Similarly, a receptive-field estimate or attention map does not establish anatomical reasoning. A stronger context test perturbs distant information while preserving local target evidence and measures whether the predicted response follows the proposed mechanism.

\begin{table*}[t]
\centering
\caption{Candidate failure mechanisms and discriminating tests. These are hypotheses to evaluate on a specified task and protocol, not family-wide causal conclusions. A positive result should be replicated and checked against alternative explanations.}
\label{tab:failure_matrix}
\footnotesize
\renewcommand{\arraystretch}{1.16}
\setlength{\tabcolsep}{4pt}
\begin{tabular}{>{\raggedright\arraybackslash}p{2.65cm}>{\raggedright\arraybackslash}p{4.55cm}>{\raggedright\arraybackslash}p{8.95cm}}
\toprule
\textbf{Mechanism} & \textbf{Candidate error signature} & \textbf{Discriminating contrast} \\
\midrule
Convolution and direct skips & Confusion on distant anatomy; nuisance detail in skip paths. & Keep backbone and data pipeline fixed. Compare direct vs. gated transfer, and separately perturb distant context vs. local artifacts. Report relational errors and boundary outcomes. \\
Dense attention and tokenization & Increasing interaction cost; errors on structures poorly represented at chosen token scale. & Vary physical token size within declared resource budget. Thin-structure errors, latency, memory; separate tokenization from decoder/initialization changes. \\
Windowed attention & Sensitivity to window placement, or restricted cross-window communication. & Translate anatomy relative to window grid. Compare trained variants with and without cross-window connections. Control padding, augmentation, physical field of view. \\
SSM / Mamba scans & Orientation sensitivity, inter-slice inconsistency, or compressed long-path errors. & Retrain matched scan variants; separately measure coordinate-transform equivariance with fixed weights. Topology, neighborhood errors, latency, memory. \\
Promptable foundation models & Domain-specific errors, prompt sensitivity, or inconsistent slice-to-slice propagation. & Keep target ontology and prompt information fixed. Compare adaptation vs. volumetric processing, then perturb prompt quality and test untouched external domains. \\
Hybrid fusion and expert routing & Branch disagreement, unstable expert use, or errors concentrated on uncertain cases. & Compare fixed fusion vs. pooled input-conditioned routing vs. spatial content routing vs. uncertainty conditioning. Calibration, utilization, errors on difficult cases, routing overhead. \\
\bottomrule
\end{tabular}
\end{table*}

\subsubsection{Bottleneck redistribution: bounded hypothesis}
\label{subsubsec:redistribution}

The \emph{redistribution framework} organizes transition-specific hypotheses. These hypotheses predict that a stated intervention produces a target benefit, with a practically relevant increase in a prespecified constraint in a declared operating regime. The contribution is to link an architectural change to both predictions and the mechanism that links them. This is not a conservation law, a no-free-lunch theorem, or a claim that every improvement comes with a cost.

For a testable instance, we need to declare the baseline, intervention, task, training conditions, target benefit, stress conditions, and finite list of protected constraints before evaluating. We need an operational measure of a constraint, like latency, boundary error, or sensitivity to a defined context perturbation, and we need a prespecified margin for a practically relevant loss. It is not enough to have a residual error that just happens to remain after an improvement, or a numerically nonzero cost below that margin, to say a bottleneck moved.

Both the intended benefit and the predicted constraint change must be measured to support a redistribution hypothesis. Conversely, a meaningful benefit together with demonstrated non-inferiority on all prespecified constraints weighs against that hypothesis in its declared scope. Introducing new constraints or stress tests after viewing the result would weaken falsifiability; such observations can instead motivate a new prospective test. This formulation permits genuine joint improvements without treating every outcome as confirmation.

\begin{table*}[t]
\centering
\caption{Six classes of transitions and the bounded questions they motivate. The source columns show architectural evidence and local experimental anchors; they do not establish that the predicted constraint is active in every model/task. Numerical comparisons stay in their original protocols.}
\label{tab:redistribution_evidence}
\scriptsize
\renewcommand{\arraystretch}{1.12}
\setlength{\tabcolsep}{3pt}
\begin{tabular}{>{\raggedright\arraybackslash}p{2.7cm}>{\raggedright\arraybackslash}p{3.05cm}>{\raggedright\arraybackslash}p{3.15cm}>{\raggedright\arraybackslash}p{3.45cm}>{\raggedright\arraybackslash}p{3.8cm}}
\toprule
\textbf{Transition} & \textbf{Evidence anchor} & \textbf{Intended benefit} & \textbf{Constraint to test} & \textbf{Bounded diagnostic question} \\
\midrule
Direct skips to aligned/routed fusion
& Skip designs in Table~\ref{table:skip_connections}; skip ablation in TransUNet \citep{chen2021transunet}
& More useful high-resolution information reaching decoder.
& Fusion overhead, nuisance transfer, or unstable routing.
& Does benefit persist when fusion cost and artifact sensitivity are measured against the same backbone? \\
CNN locality to CNN--transformer hybridization
& Encoder/decoder contrasts in TransUNet \citep{chen2021transunet}; Table~\ref{tab:transformer_medseg}
& More direct interaction of spatially separated features.
& Tokenization, interaction cost, or branch alignment.
& Does gain on controlled relational cases persist with matched initialization and deployment budgets? \\
Dense to windowed interaction
& Shifted-window ablation in Swin \citep{review95_liu2021swin}; medical designs in Table~\ref{tab:transformer_medseg}
& Higher spatial resolution interaction within resource budget.
& Dependence on window connectivity and placement.
& Does observed long-span error respond specifically to cross-window connections? \\
Attention to selective state propagation
& Recurrence and scan designs in Table~\ref{tab:mamba_seg}
& Long-range processing without dense attention matrix.
& Ordering, state compression, realized kernel cost.
& Do trained scan variants change predicted topology/orientation error at comparable resources? \\
Generic prompting to medical adaptation
& Adaptation mechanisms in Table~\ref{tab:detailed_comparison}
& Better alignment with target domain or volumetric task.
& Prompt information, adaptation cost, transfer outside adaptation distribution.
& Does improvement persist with equivalent prompts, untouched external evaluation? \\
Fixed fusion to conditional expert allocation
& TA-MoSC paper and released implementation \citep{luo2025tamosc_verified,luo2025tamosc_code_verified}
& Input-dependent selection of feature paths.
& Router calibration, utilization, stability, overhead.
& Does conditioning add benefit over strong input-conditioned baseline at declared inference budget? \\
\bottomrule
\end{tabular}
\end{table*}

\paragraph{Three worked contrasts and their limits}
\emph{TransUNet token granularity.} In the original analytical study, reducing patch size from 16 to 8 increases the sequence length from 196 to 784 and the reported Synapse DSC from 77.48 to 77.83~\citep{chen2021transunet}. At fixed attention width, this produces 16 times as many token pairs, a property of dense attention rather than a measured latency ratio. The comparison motivates a resolution--interaction cost hypothesis but does not establish the local-precision mechanism or a realized computational penalty without matched boundary, topology, latency, and memory measurements.

\emph{Swin cross-window communication.} In the original Swin-T ablation, regular and shifted windows achieve ADE20K mIoU values of 43.3 and 46.1, respectively~\citep{review95_liu2021swin}. This supports the relevance of the connectivity pattern in that general-vision setting, but it is neither a medical segmentation result nor proof of a particular anatomical error. The window-placement and long-range relational tests in Table~\ref{tab:failure_matrix} would be needed to establish the proposed medical mechanism. Restricted attention alone does not imply a measurable penalty.

\emph{Swin-UMamba initialization.} The within-architecture comparison reports AbdomenMRI DSC of 0.7054 without ImageNet pretraining and 0.7760 with pretraining, with unchanged reported parameter and inference-FLOP counts~\citep{liu2024swin}. Thus, the reported segmentation accuracy improves without an increase in the stated inference cost when initialization changes. Pretraining resources belong in a life-cycle cost analysis only when that scope is specified in advance; they cannot be introduced retrospectively to invalidate an inference-only comparison. This contrast supports an initialization effect in its reported setting, not a universal advantage or a claim of external-domain robustness.

\paragraph{Evidence status across transitions}
The six transitions have different levels of support in the records reviewed here. The skip-path, CNN--transformer, and medical prompt-adaptation literature contains source-specific performance comparisons, but not all of the proposed secondary constraints are measured under matched conditions. The shifted-window example isolates connectivity in general vision rather than medical anatomy, whereas the Swin-UMamba example isolates initialization rather than a scan-induced bottleneck. SSM scan design and conditional routing are documented architecturally, but decisive comparisons of topology, calibration, and realized cost remain limited in the audited evidence. Accordingly, these records provide mechanism anchors and selected local contrasts, not a complete demonstration of constraint redistribution across all six transitions. This evidence gap motivates the discriminating tests in Table~\ref{tab:research_roadmap}; it is not a negative performance judgment about any method family.

\paragraph{From illustrative contrast to test}
Mechanism studies should report their patient splits, input geometry, label definitions, initialization, and evaluation code. Small interventions can be compared within a common training recipe, while substantially different architectures also require comparable, disclosed tuning budgets: identical optimizers and epoch counts do not guarantee equivalent search effort or training cost. Training, inference, and reused pretraining resources should be reported separately. The controlled re-evaluation by~\citet{isensee2024nnunetrevisited} reinforces the importance of strong, properly configured baselines.

Replication should use independent patients and more than one relevant operating condition. Repeated training seeds estimate optimization variability but do not make correlated image slices independent test cases. A bounded comparison can support, contradict, or remain inconclusive for the stated mechanism; it does not rank entire architectural families. Section~\ref{subsec:roadmap} provides an uncertainty-aware decision rule.
\subsection{Protocol-Stratified Evidence Admissibility Audit}
\label{subsec:tradeoff_decomp}

Comparability depends on experimental records, not architecture names. Two models evaluated at $224\times224$ pixels may still differ in patient partitions, physical field of view, class averaging, metric implementation, prompts, initialization, or model scope. Table~\ref{tab:tradeoff_quantitative} specifies the admissible interpretation of representative Synapse records: Tier A denotes verified comparisons within an evaluation, Tier B retains related but incompletely matched records, and Tier C covers incompatible dimensionality, task, label space, or inference conditions.

Experimental provenance is distinct from architectural authorship. An independent implementation of U-Net provides primary evidence for its own evaluation, even though the architecture was introduced elsewhere; a subsequent transcription of that result is secondary reporting. Accordingly, the medical reimplementations in~\citet{azad2024unet_verified} are attributed to the evaluation that generated the numerical values, not solely to the original architectural publication.

\begin{table*}[t]
\centering
\caption{Provenance and admissibility of representative Synapse records. Tiers refer to the scope of comparison stated. HD is kept if it's the label of the source. Code-level percentile alone doesn't mean physical units, aggregation, exact implementation of a published checkpoint. No cross-paper leaderboard is built.}
\label{tab:tradeoff_quantitative}
\scriptsize
\renewcommand{\arraystretch}{1.16}
\setlength{\tabcolsep}{3pt}
\begin{tabular}{>{\raggedright\arraybackslash}p{2.5cm}>{\raggedright\arraybackslash}p{2.9cm}>{\raggedright\arraybackslash}p{3.6cm}>{\raggedright\arraybackslash}p{3.0cm}>{\raggedright\arraybackslash}p{4.15cm}}
\toprule
\textbf{Experimental record} & \textbf{Source and geometry} & \textbf{Boundary metric provenance} & \textbf{Admissible scope} & \textbf{Remaining restriction} \\
\midrule
U-Net, UNet++ and Residual U-Net reimplementations
& Azad et al.; 2D $224\times224$ \citep{azad2024unet_verified}
& Author preprint Table II(c) HD for boundary measure
& A inside that evaluation. B to separately trained papers
& Primary reimplementation results. These Synapse values, HD95 semantics are not established in original model publications. \\
TransUNet and its ablations
& Chen et al.; 2D $224\times224$ (unless ablation changes geometry) \citep{chen2021transunet}
& Paper: HD. Released \texttt{utils.py} \texttt{hd95} (no voxel spacing arg) \citep{transunet_code_verified}
& A stated within paper contrasts. B unverified cross-paper matches
& Keep paper label. Disclose code difference. Verification of physical units, checkpoint-to-code correspondence required. \\
VM-UNet
& Original 2D Synapse evaluation \citep{ruan2024vm}
& Original paper HD95
& B for comparison to other papers listed here
& Metric name and input dimensions don't verify the whole protocol tuple. Don't transfer costs from another input configuration. \\
Swin-Unet
& Original Synapse evaluation \citep{cao2022swinunet}
& Keep paper HD label, need explicit tracing of implementation semantics
& B cross papers
& Don't reject or admit record just based on HD/HD95 abbreviation. Still relevant partition, aggregation, code provenance. \\
PVT-EMCAD-B2
& Original Synapse evaluation \citep{rahman2024emcad}
& Keep source specific boundary metric
& B cross papers
& Don't establish full comparison protocol, cost scope from same dataset, nominal resolution. \\
SAMed
& Automatic semantic segmentation of Synapse \citep{zhang2023samed}
& Keep source terminology, until we link to evaluation implementation
& B for potentially comparable automatic settings
& Verify inference and conditioning before comparison. SAM lineage alone doesn't mean manual test prompts. \\
nnU-Net later comparison
& Volumetric configuration, trace to experiment generating value \citep{isensee2021nnu}
& Source specific metric and geometry
& C relative to this 2D stratum
& Original nnU-Net citation is architectural provenance, unless they report the exact experiment. \\
nnFormer, UNETR++
& Original volumetric evaluations \citep{zhou2023nnformer,shaker2024unetrpp}
& Original source specific boundary measures
& C relative to this 2D stratum
& Keep as volumetric evidence. Only compare numerically in separately verified 3D protocol. \\
\bottomrule
\end{tabular}
\end{table*}

Differences between papers and released evaluation code can affect metric interpretation. The TransUNet preprint labels a boundary measure HD, while the released evaluation utility calls a 95th-percentile Hausdorff routine without passing voxel spacing~\citep{chen2021transunet,transunet_code_verified}. This clarifies the inspected implementation but does not justify silently relabeling every published score or assuming physical units. Both descriptions should remain in the source audit until checkpoint and protocol correspondence are established.

Efficiency comparisons require a separate accounting record covering included model components, input tensor geometry, batch size, numerical precision, hardware or FLOP counter, and pre/post-processing. Decoder-only and whole-pipeline costs answer different questions. Missing standard deviations, confidence intervals, or protocol fields cannot be reconstructed from reported point estimates; the heterogeneous entries here therefore do not support cross-study significance tests, rank correlations, or empirical Pareto-frontier claims. Cross-domain robustness should be assessed on held-out domains sharing a task and label ontology.

\subsection{Open Challenge 1: Native 3D State Space Modeling}
\label{subsec:native3d}

A key design requirement for volumetric SSMs is to specify how spatial neighbors exchange information. SegMamba, nnMamba, and LKM-UNet use different directional, channel--spatial, and multiscale constructions~\citep{xing2024segmamba,gong2025nnmamba,wang2024lkm}. Their mechanisms should be compared using measured neighborhood preservation and task performance rather than the designation \emph{native 3D} alone.

In a single linear sequence, each interior token has at most two immediate sequence neighbors, whereas an interior voxel in a 3D grid has six face-sharing neighbors. A one-to-one ordering therefore cannot make every grid edge an adjacent sequence pair. This combinatorial limitation does not imply segmentation failure: multiple scan paths, local convolutions, and recurrent states may preserve the information needed for a given task. It does motivate measuring the relationships among scan-path length, orientation, state compression, and anatomical error.

Three directions merit controlled evaluation. \emph{Spatial graph propagation} retains explicit neighborhood edges, although its efficiency depends on sparsity, memory access, and state-update implementation. \emph{Learned ordering} can adapt propagation paths to the input, but the cost of constructing that ordering must be included. \emph{Local or hierarchical ordering} provides a more constrained alternative and can be studied through path-distortion and topology analyses. Separately trained scan variants should be compared under similar budgets, with coordinate-transform tests and measured latency, peak memory, overlap, and surface outcomes. None of these strategies should be presumed efficient solely from its representation.

\subsection{Open Challenge 2: Large-Scale Medical Foundation Pretraining}
\label{subsec:pretraining}

The Swin-UMamba contrast motivates separating initialization effects from architectural effects. A controlled study could cross backbone choice with random, natural-image, and medical-image initialization while documenting pretraining data and resources, downstream tuning, label efficiency, and held-out-domain transfer. Improvement on one of these endpoints would not establish improvement on the others.

A medical pretraining corpus should be designed for its intended use. Volumetric applications require information about slice spacing and spatial correspondence; multisite use additionally requires patient, scanner, and pathology diversity. Transfer claims need documented separation between pretraining subjects and final evaluation cohorts. A single model need not cover every modality, but its coverage, exclusions, and access conditions should be explicit enough to audit the claimed scope.

Masked reconstruction may encourage anatomically useful representations, but the pretraining objective alone does not establish preservation of the relations needed for segmentation. Complementary objectives should therefore be evaluated using downstream probes and controlled interventions. \emph{Cross-view consistency} can align corresponding anatomical features in reformatted views, while distinguishing invariant global identity from spatially equivariant dense representations; unaligned axial and coronal sections need not contain the same structures. \emph{Deformation equivariance} tests whether predictions or feature fields transform consistently under plausible known deformations. \emph{Anatomical relation prediction} can assess relative position, connectivity, or containment, provided that pathological variation and uncertainty are not treated as errors by construction.

The decisive evaluation should hold the downstream architecture and data budget fixed, report performance across labeled-data fractions, and reserve independent institutions or acquisition domains for final testing. Medical pretraining should be compared with strong reconstruction and contrastive baselines, with protected endpoints for unusual pathology and small structures. Excessively restrictive anatomical priors may otherwise improve average consistency while obscuring clinically important variation.

\subsection{Open Challenge 3: Adaptive Feature Fusion}
\label{subsec:adaptive_fusion}

Dynamic U-Net investigates adaptive feature calibration for abdominal multi-organ segmentation~\citep{yang2025dynamic}. TA-MoSC provides a related example of input-conditioned expert selection~\citep{luo2025tamosc_verified,luo2025tamosc_code_verified}: its released implementation pools features globally to produce stage-specific top-$k$ expert assignments. The router parameters remain fixed at inference, but its decisions depend on the input. Spatially resolved uncertainty conditioning is therefore a possible extension, whose incremental value should be tested against existing pooled and spatial content-conditioned gates.

\paragraph{Spatial alignment and conditioning}
Let $\mathbf{E}_j(x)$ denote encoder features at level $j$, and let $\Pi_{ij}$ align their spatial resolution and channel dimension with decoder stage $i$. A candidate fusion rule is
\begin{equation}
  \mathbf{O}_i(x,p)=\sum_{j=1}^{L}g_{ij}(x,p)
  \bigl[\Pi_{ij}(\mathbf{E}_j(x))\bigr]_p,
\end{equation}
where $p$ is spatial position. $g_{ij}\geq0$. $\sum_jg_{ij}=1$. So every term in the sum is a feature vector in the same decoder space.

Let $s_{ij}(x,p)$ be a content-compatibility score and $\hat{\sigma}_i^2(x,p)$ a pre-gate uncertainty estimate, for example from an auxiliary head on the coarser decoder state. This estimate must be available without reference labels at inference. After a specified calibration step, a candidate gate is
\begin{equation}
  g_{ij}(x,p)=
  \frac{\exp\!\left(s_{ij}(x,p)+\lambda_j\hat{\sigma}_i^2(x,p)\right)}
       {\sum_{k=1}^{L}\exp\!\left(s_{ik}(x,p)+\lambda_k\hat{\sigma}_i^2(x,p)\right)}.
\end{equation}
Only differences among the $\lambda_j$ coefficients alter the relative weights: an identical additive uncertainty term cancels under softmax. High uncertainty does not necessarily favor shallow features, since it may also reflect noise or domain shift. The signs and magnitudes of the coefficients, and any benefit beyond content-only routing, therefore require empirical evaluation.

A decisive comparison should include fixed fusion, pooled input-conditioned routing, spatial content-only routing, and uncertainty-conditioned routing under matched backbone, supervision, and information budgets. Inference accounting must include auxiliary uncertainty heads. Calibration, expert utilization, difficult-case errors, performance under domain shift, and routing overhead should be reported; nonuniform expert use alone does not establish useful specialization.

\paragraph{Information gain as a conceptual diagnostic}
A complementary question is whether an added feature contributes target information beyond the decoder state and features already selected. Let $\widetilde{\mathbf E}_{ij}=\Pi_{ij}(\mathbf E_j)$ denote an aligned feature and $\mathcal S_i$ the selected levels. A conceptual selection criterion is
\begin{equation}
\label{eq:conditional_mi}
j^*=\arg\max_{j\notin\mathcal S_i}
I\!\left(Y;\widetilde{\mathbf E}_{ij}\mid
\mathbf D_{i+1},\{\widetilde{\mathbf E}_{ik}\}_{k\in\mathcal S_i}\right).
\end{equation}
Here $\mathbf D_{i+1}$ is the coarser decoder state, and $Y$ is the target variable in the analysis distribution rather than an annotation available to the inference-time router. Estimating conditional information in high-dimensional features is difficult, and greedy selection does not guarantee an optimal feature subset. Equation~(\ref{eq:conditional_mi}) therefore specifies a diagnostic research question, not an implemented algorithm or a validated segmentation objective.

\subsection{Open Challenge 4: Frequency-Domain Feature Learning for Boundary Reconstruction}
\label{subsec:frequency}

Spectral-bias theory concerns how neural networks fit frequency components of functions defined over the input domain~\citep{jacot2018ntk,rahaman2019spectral,basri2019convergence,xu2020frequency}. These frequencies are not identical to the spatial Fourier components of an output segmentation mask. Connecting the two requires assumptions about the network, data distribution, loss, and spatial representation. The cited theory therefore motivates a hypothesis; it does not establish that medical-image boundaries must be learned later than region interiors.

A slower improvement in HD95 than DSC can result from rare surface errors, class imbalance, annotation noise, or differing metric sensitivity and is not, by itself, evidence of spectral bias. A discriminating study should measure spatial spectral residuals on fixed validation cases, separate thin structures from uncertain annotations, and test whether a frequency-specific intervention improves prespecified endpoints beyond an equally strong spatial boundary objective.

Frequency-aware segmentation already has established precedents. T-Mamba and HybridMamba integrate spectral processing~\citep{hao2026t,wu2025hybridmamba}; EM-Net integrates frequency learning with 3D Mamba~\citep{verified_chang2024emnet}; FaRMamba evaluates wavelet, cosine, and Fourier variants~\citep{verified_rong2025farmamba}; and WMC-Net and FreqConvMamba create wavelet or mixed spectral pathways~\citep{verified_fan2026wmcnet,verified_song2026freqconvmamba}. Wavelet fusion also precedes these SSM architectures~\citep{zhou2023xnet}. These precedents limit the open question to when a spectral pathway enhances boundary fidelity, topology, or robustness under a controlled protocol. Overlap gains alone do not resolve this.

\paragraph{A controlled loss and pathway comparison}
An example of an auxiliary objective is the unitary discrete Fourier transform $\mathcal F$ of predicted class probabilities $\hat P_c$ and reference masks $Y_c$:
\begin{equation}
\begin{aligned}
\mathcal L &= \mathcal L_{\mathrm{seg}}+\eta\mathcal L_{\mathrm{freq}},\\
\mathcal L_{\mathrm{freq}}
 &=\sum_{c,\omega}w(\omega)
 \left|\mathcal F[\hat P_c](\omega)-\mathcal F[Y_c](\omega)\right|^2 .
\end{aligned}
\end{equation}
The transform normalization, class weights, physical sampling, and $\eta\geq0$ must be specified. A choice $w(\omega)\propto\|\omega\|^\alpha$, $\alpha>0$, emphasizes high spatial frequencies but assigns no weight to the DC component; the ordinary segmentation term maintains supervision of region content. High-frequency emphasis can also amplify artifacts and annotation noise, so a benefit is not guaranteed.

A factorial experiment should evaluate a baseline, a spatial boundary loss, a frequency pathway, and their combination. Wavelet, Fourier, and cosine variants can then be evaluated within the same backbone and tuning budget. Predefine surface and thin-structure endpoints, spectral residuals, topology, and robustness to noise and resampling; report latency and memory for the added transforms. This design evaluates whether a spectral pathway adds information beyond stronger spatial supervision. An accompanying signal-processing analysis can examine what pooling, tokenization, and reconstruction preserve, without assuming that every boundary error is a frequency-learning failure.
\begin{table*}[!htbp]
\centering
\caption{Research roadmap with explicit uncertainty-aware decisions. Each row uses the benefit and protected-loss contrasts defined in the text. The proposed comparisons should share evaluation data and code and use reported, comparable tuning/search budgets. A contradicted cost-qualified claim does not imply that the underlying mechanism has no benefit in another regime.}
\label{tab:research_roadmap}
\scriptsize
\renewcommand{\arraystretch}{1.13}
\setlength{\tabcolsep}{3pt}
\begin{tabular}{>{\raggedright\arraybackslash}p{2.35cm}>{\raggedright\arraybackslash}p{3.25cm}>{\raggedright\arraybackslash}p{5.05cm}>{\raggedright\arraybackslash}p{4.0cm}>{\raggedright\arraybackslash}p{1.5cm}}
\toprule
\textbf{Direction} & \textbf{Bounded hypothesis} & \textbf{Controlled experiment} & \textbf{Evidence against the stated claim} & \textbf{Constructs informed} \\
\midrule
Volumetric SSM propagation
& A specified topology-aware path improves a topology or orientation endpoint without unacceptable deployment cost.
& Retrain raster, multidirectional, and proposed paths in the same backbone with comparable tuning and training budgets. Test multiple relevant volumetric settings; report topology, surface error, equivariance, latency, and memory.
& $U_\Delta<\delta_{\mathrm{3D}}$ relative to the strongest matched path, or $L_r>m_r$ for a protected cost. Fixed-weight sensitivity alone is insufficient.
& LP, GC, CE \\
\midrule
Medical pretraining
& Medical initialization improves a specified external-domain or low-label endpoint beyond matched alternative initialization.
& Cross architecture with initialization; record pretraining scale and resources separately. Match downstream data and tuning budgets, prevent patient overlap, and reserve external domains for final evaluation.
& $U_\Delta<\delta_{\mathrm{pre}}$, or a protected pathology, domain, or resource loss satisfies $L_r>m_r$.
& CG; LP/GC with dedicated probes \\
\midrule
Spatial uncertainty-conditioned fusion
& Pre-gate uncertainty adds useful routing information beyond spatial content and pooled input-conditioned gates.
& Compare fixed, pooled input-conditioned, spatial content-only, and uncertainty-conditioned fusion. Match available supervision and budget; include uncertainty-head cost. Measure calibration, utilization, and difficult-case errors.
& $U_\Delta<\delta_{\mathrm{route}}$ against the strongest matched gate, or a protected calibration/robustness/cost loss has $L_r>m_r$.
& LP, CG; GC if tested \\
\midrule
Frequency-aware reconstruction
& A specified spectral pathway improves surface or thin-structure accuracy beyond stronger spatial supervision.
& Factorially compare boundary loss and frequency pathway; compare transforms within one backbone. Control geometry, augmentation, and tuning. Track spectral residuals, topology, noise sensitivity, and cost.
& $U_\Delta<\delta_{\mathrm{freq}}$ against the boundary-loss baseline, or a protected noise/shift/cost loss has $L_r>m_r$.
& LP, CE; CG if tested \\
\midrule
Multi-objective selection
& Equal-budget selection produces a useful operating-point improvement on a fixed outcome set.
& Compare single- and multi-objective procedures using development domains and equal search budgets. Freeze candidates before external evaluation. Prespecify the reference set, endpoints, margins, and frontier statistic.
& The adjusted upper bound on the prespecified frontier improvement is below its worthwhile margin, or a protected loss has $L_r>m_r$.
& CE, CG and measured task outcomes \\
\bottomrule
\end{tabular}
\end{table*}

\subsection{Open Challenge 5: Multi-Objective Architectural Optimization}
\label{subsec:moo}

Multi-objective architecture search for medical segmentation has established precedents, including CTU-NAS~\citep{ying2023ctunas_verified}. The question developed here concerns explicit boundary, deployment, and held-out-domain constraints; it does not claim to introduce multi-objective search. Architecture search selects a design and should be distinguished from input-conditioned routing during inference.

The LP--GC--CE--CG framework organizes capabilities and evaluation questions. A practical optimization problem must instead use specified observable endpoints. Boundary error and overlap are useful outcomes, but they do not by themselves validate local-detail preservation or long-range reasoning. The following formulation is an illustrative outcome-based design problem, not a four-dimensional measurement model of the latent constructs.

For a fixed task and evaluation protocol, let $H(A)$ be a specified surface-error measure and $D(A)$ an overlap score. Optional display coordinates are
\begin{equation}
B_{\mathrm{perf}}(A)=1-\frac{H(A)-H_{\min}}{H_{\max}-H_{\min}},
\end{equation}
and
\begin{equation}
O_{\mathrm{perf}}(A)=\frac{D(A)-D_{\min}}{D_{\max}-D_{\min}}.
\end{equation}
The bounds are fixed from development data before comparing candidates and have nonzero ranges. Raw outcomes and confidence intervals remain the primary report. These monotone transformations preserve ordering; they add no evidence and need not remain between zero and one for a new model outside the calibration range.

For one explicitly chosen deployment cost, such as median whole-pipeline latency $F(A)>0$, define
\begin{equation}
CE_{\mathrm{lat}}(A)
=1-\frac{\log F(A)-\log F_{\min}}
{\log F_{\max}-\log F_{\min}},
\label{eq:ce_perf}
\end{equation}
with $0<F_{\min}<F_{\max}$. Input geometry, batch size, precision, hardware, warm-up, and pre/post-processing are fixed. Memory and energy are additional objectives or protected constraints, not silently absorbed into latency. Parameter count and FLOPs can describe the model but do not replace this measurement.

For domains $\mathcal D^{(t)}$ sharing clinical task $t$, label ontology, and aggregation, an external overlap summary can be written as
\begin{equation}
G(A;\mathcal D^{(t)})=
\sum_{d\in\mathcal D^{(t)}}w_d\,\frac{\mathrm{DSC}(A,d)}{100},
\qquad \sum_d w_d=1.
\end{equation}
Domain weights are nonnegative and prespecified, and per-domain results are also reported so that an average cannot conceal a severe failure. Incompatible tasks are kept separate. This score measures external overlap, not every aspect of generalization; calibration and worst-domain behavior may require separate endpoints.

Architecture selection can then seek
\begin{equation}
\underset{A\in\mathcal A}{\operatorname{ParetoMax}}
\bigl(B_{\mathrm{perf}},O_{\mathrm{perf}},CE_{\mathrm{lat}},G_{\mathrm{dev}}\bigr)(A),
\end{equation}
subject to declared deployment and clinical-performance constraints. Search and tuning use development domains only. Candidates are frozen before evaluation on untouched external domains, where $G_{\mathrm{ext}}$ and the other outcomes are reported. Repeatedly selecting on those external results would invalidate their held-out role. Dedicated relational, topology, and perturbation tests must be added whenever conclusions about GC or LP are sought.

The research question is whether multi-objective selection produces better operating choices than strong single-objective and existing multi-objective baselines under equal search budgets. Adaptive computation is one possible design variable, not a requirement. Its total cost includes gating, expert dispatch, and variability across cases; mean latency alone can hide a clinically relevant tail. No empirical Pareto frontier is inferred from the heterogeneous published values in this review.

\subsection{A Unified Architectural Roadmap}
\label{subsec:roadmap}

Table~\ref{tab:research_roadmap} translates the preceding questions into bounded tests. A successful direction need not improve every outcome. It must demonstrate the specified benefit while meeting the constraints attached to that claim. The experiments can be conducted independently, and their usefulness does not depend on assembling all five mechanisms in one model.


\paragraph{Common decision rule}
Before comparison, specify the primary endpoint, minimum worthwhile benefit $\delta>0$, protected endpoints and tolerated losses $m_r\geq0$, patient-level analysis, training repetitions, confidence level, and multiplicity control. Orient the target contrast $\Delta$ so that positive values favor the intervention, and each protected loss $\Gamma_r$ so that positive values are harmful. Write their suitably adjusted confidence intervals as $[L_\Delta,U_\Delta]$ and $[L_r,U_r]$.

Evidence for the claimed benefit requires $L_\Delta>\delta$; evidence against a benefit of that size requires $U_\Delta<\delta$. A protected endpoint meets its non-inferiority criterion when $U_r<m_r$, whereas $L_r>m_r$ establishes an unacceptable loss. Other cases are inconclusive for that criterion. A point estimate below a margin, failure to reach significance, or an interval crossing a margin is not automatically evidence against the mechanism. These rules apply to the prespecified benefit-and-constraint claim, not to the value of an entire research area.

For a redistribution test, the targeted benefit and the predicted constraint change must both be observed with adequate precision. If the target improves and every prespecified constraint meets non-inferiority, the proposed harmful redistribution is contradicted within that scope. If both benefit and a relevant constraint loss are established, the result supports the proposed shift and quantifies its practical cost. Neither outcome establishes a universal property of architectural progress.

The roadmap is anchored to the mechanisms in Tables~\ref{tab:failure_matrix} and~\ref{tab:redistribution_evidence}. It does not estimate how often failures occur in the literature. Its practical output is a set of comparisons that distinguish useful mechanisms from gains attributable to initialization, supervision, resources, or evaluation changes.

\subsection{Review Limitations and Threats to Validity}
\label{subsec:limitations}

The purposive narrative design supports mechanism coverage rather than field-wide prevalence estimates. Publication and survivorship bias favor successful variants; negative results, routing collapse, unstable optimization, and unsuccessful adaptations are less visible. Within-paper ablations provide stronger local evidence than comparisons of separate papers, but they still depend on the authors' experimental choices and do not establish a family-wide effect.

Several protocol details remain incompletely documented in the source literature. Code can clarify an implementation while leaving the exact published checkpoint, data preprocessing, or metric aggregation unresolved. Source-specific point estimates are therefore retained with their provenance rather than pooled. No new training experiment or independent replication is reported here.

The LP and GC constructs also require more validation than a metric table supplies. Boundary and overlap outcomes, relational probes, topology tests, and external-domain results offer complementary evidence; none alone proves a general capability. The proposed transition hypotheses are bounded by their tasks, perturbations, resources, and measured endpoints. Broadening that scope requires new evidence.

\section{Conclusion}

Architectural progress in medical image segmentation is best assessed through the intervention, task, and operating conditions that produced a reported result. Convolution, attention, selective state propagation, pretrained prompting, and conditional fusion offer different ways to construct and exchange representations; the architectural family name alone does not establish suitability for a particular clinical task.

This review evaluates these mechanisms in relation to local precision, global context, computational efficiency, and cross-domain generalization. The bounded redistribution framework links each intended benefit to a separately testable constraint. A demonstrable benefit accompanied by non-inferiority on every prespecified protected outcome does not support the proposed harmful constraint shift within that scope. The three worked contrasts illustrate why mechanism attribution requires more than a benchmark improvement or an unmatched cost comparison.

Three reporting practices follow from the evidence audit. First, distinguish the architectural origin from the experiment that generated each reported number, and retain unresolved protocol fields. Second, report overlap and boundary outcomes with patient-level uncertainty while using dedicated tests for topology and long-range reasoning. Third, measure deployment costs under fixed accounting conditions and reserve external domains for evaluation after model selection.

The resulting research agenda includes volumetric propagation, medical pretraining, uncertainty-conditioned fusion, spectral reconstruction, and multi-objective selection. These are alternatives to compare against strong baselines, not mandatory components of a larger model. Progress is demonstrated by reproducible, practically meaningful gains under declared constraints.

\section*{Declaration of competing interest}
The authors state that there are no known competing financial interests or personal relationships that may appear to have influenced the work presented in this paper.

\section*{CRediT authorship contribution statement}
 \textbf{Nand Kumar Yadav:} Conceptualization, Methodology, Investigation, Writing -- original draft. \textbf{Smriti Regmi:} Investigation, Data curation. \textbf{Moinak Bose:} Investigation, Visualization. \textbf{N. Reed Alexander:} Writing -- review \& editing. \textbf{Debesh Jha:} Supervision, Writing -- review \& editing. \textbf{William CW Chen:} Supervision, Writing -- review \& editing.

\section*{Data availability}
No patient-level data or new experimental measurements were generated. Supplementary Data S1 (\path{method_register.csv}) provides the method inventory and family assignments. Supplementary Data S2 (\path{numerical_verification.csv}) provides the source-verification records, inspected configurations, and unresolved fields. Both files accompany this manuscript, together with the cell-level revision audit, the contextual-publication index, and the explanatory Supplementary Information. Published datasets and implementations are identified through their cited sources.

\section*{Declaration of generative AI and AI-assisted technologies in the manuscript preparation process}
During manuscript preparation, ChatGPT (OpenAI) assisted with language editing, literature organization, drafting of source-verification records, mathematical exposition, and schematic drawing code and layout. GPT-6 Astra Pro was used for the final revision; the historical model versions used for earlier figure preparation were not recorded in the project files. Assistance with individual figures is identified in their captions. The authors retain full responsibility for source verification, interpretation, originality, and approval of the final submitted content.
\section*{Acknowledgements}
This work received partial support from the National Science Foundation (Grant No. \href{https://www.nsf.gov/awardsearch/showAward?AWD_ID=2346643}{2346643}), the U.S. Department of Defense (Award No. \href{https://dtic.dimensions.ai/details/grant/grant.14525543}{FA9550-23-1-0495}), and the U.S. Department of Education (Grant No. P116Z240151). Any opinions, interpretations, conclusions, or recommendations expressed herein are those of the authors and should not be construed as representing the official views of the funding agencies.


\bibliographystyle{elsarticle-harv-en}
\bibliography{referencef,review_corrections,additional_review_refs,evidence_implementation,numeric_sources,new_papers_verified,transformer_numeric_refs}


\end{document}